\Chapter{37}

\Verse{1}
{
    \zA{话说史湘云回家后，宝玉等仍不过在园中嬉游吟咏不提。}
}
{
    \eA{But to continue. After Shih Hsiang-yün's return home, Pao-yü and the other
        inmates spent their time, as of old, in rambling about in the garden in
        search of pleasure,}
}
{
}

\Verse{2}
{
    \zA{且说贾政自元妃归省之后，居官更加勤慎，以期仰答皇恩。}
}
{
    \eA{and in humming poetical compositions. But without further reference to their
        doings, let us take up our narrative with Chia Cheng. Ever since the visit
        paid to her home by the imperial consort,}
}
{
}

\Verse{3}
{
    \zA{皇上见他人品端方， 风声清肃，虽非科第出身，却是书香世代，因特将他点了学差，也无非是选拔真才 之意。}
}
{
    \eA{he fulfilled his official duties with additional zeal, for the purpose of
        reverently making requital for the grace shown him by the Emperor. His
        correct bearing and his spotless reputation did not escape His Majesty's
        notice, and he conferred upon him the special appointment of Literary
        Chancellor,}
}
{
}

\Verse{4}
{
    \zA{这贾政只得奉了旨，择于八月二十日起身。}
}
{
    \eA{with the sole object of singling out his true merit; for though he had not
        commenced his career through the arena of public examinations, he belonged
        nevertheless to a family addicted to letters during successive generations.}
}
{
}

\Verse{5}
{
    \zA{是日拜别过宗祠及贾母，便起身 而去。}
}
{
    \eA{Chia Cheng had, therefore, on the receipt of the imperial decree, to select
        the twentieth day of the eighth moon to set out on his journey.}
}
{
}

\Verse{6}
{
    \zA{宝玉等如何送行，以及贾政出差外面诸事，不及细述。}
}
{
    \eA{When the appointed day came, he worshipped at the shrines of his ancestors,
        took leave of them and of dowager lady Chia, and started for his post.}
}
{
}

\Verse{7}
{
    \zA{单表宝玉自贾政起身之后，每日在园中任意纵性游荡，真把光阴虚度，岁月空 添。}
}
{
    \eA{It would be a needless task, however, to recount with any full particulars
        how Pao-yü and all the inmates saw him off, how Chia Cheng went to take up
        his official duties, and what occurred abroad, suffice it for us to notice
        that Pao-yü,}
}
{
}

\Verse{8}
{
    \zA{这日甚觉无聊，便往贾母王夫人处来混了一混，仍旧进园来了。}
}
{
    \eA{ever since Chia Cheng's departure, indulged his caprices, allowed his
        feelings to run riot, and gadded wildly about. In fact, he wasted his time,
        and added fruitless days and months to his age.}
}
{
}

\Verse{9}
{
    \zA{刚换了衣裳， 只见翠墨进来，手里拿着一幅花笺，送与他看。}
}
{
    \eA{On this special occasion, he experienced more than ever a sense of his lack
        of resources, and came to look up his grandmother Chia and Madame Wang. With
        them, he whiled away some of his time,}
}
{
}

\Verse{10}
{
    \zA{宝玉因道：“可是我忘了，才要瞧 瞧三妹妹去。}
}
{
    \eA{after which he returned into the garden. As soon as he changed his costume,
        he perceived Ts'ui Mo enter, with a couple of sheets of fancy notepaper,}
}
{
}

\Verse{11}
{
    \zA{你来的正好。}
}
{
    \eA{in her hand, which she delivered to him.}
}
{
}

\Verse{12}
{
    \zA{可好些了？”}
}
{
    \eA{"It quite slipped from my mind,"}
}
{
}

\Verse{13}
{
    \zA{翠墨道：“姑娘好了，今儿也不吃药了， 不过是冷着一点儿。”}
}
{
    \eA{Pao-yü remarked. "I meant to have gone and seen my cousin Tertia; is she
        better that you come?" "Miss is all right," Ts'ui Mo answered. "She hasn't
        even had any medicine to-day. It's only a slight chill."}
}
{
}

\Verse{14}
{
    \zA{宝玉听说，便展开花笺看时，上面写道： 妹探谨启二兄文几：前夕新霁，月色如洗，因惜清景难逢，未忍就卧，漏已三
        转，犹徘徊桐槛之下，竟为风露所欺，致获采薪之患。}
}
{
    \eA{When Pao-yü heard this reply, he unfolded the fancy notepaper. On perusal,
        he found the contents to be: "Your cousin, T'an Ch'un, respectfully lays
        this on her cousin Secundus' study-table. When the other night the blue sky
        newly opened out to view, the moon shone as if it had been washed clean!
        Such admiration did this pure and rare panorama evoke in me that I could not
        reconcile myself to the idea of going to bed.}
}
{
}

\Verse{15}
{
    \zA{昨亲劳抚嘱已，复遣侍儿问 切，兼以鲜荔并真卿墨迹见赐，抑何惠爱之深耶!今因伏几处默，忽思历来古人，
        处名攻利夺之场，犹置些山滴水之区，远招近揖，投辖攀辕，务结二三同志，盘桓 其中，或竖词坛，或开吟社：虽因一时之偶兴，每成千古之佳谈。}
}
{
    \eA{The clepsydra had already accomplished three turns, and yet I roamed by the
        railing under the dryandra trees. But such poor treatment did I receive from
        wind and dew (that I caught a chill), which brought about an ailment as
        severe as that which prevented the man of old from picking up sticks. You
        took the trouble yesterday to come in person and cheer me up. Time after
        time also did you send your attendants round to make affectionate inquiries
        about me. You likewise presented me with fresh lichees and relics of
        writings of Chen Ch'ing.}
}
{
}

\Verse{16}
{
    \zA{妹虽不才，幸叨 陪泉石之间，兼慕薛林雅调。}
}
{
    \eA{How deep is really your gracious love! As I leant to-day on my table plunged
        in silence, I suddenly remembered that the ancients of successive ages were
        placed in circumstances,}
}
{
}

\Verse{17}
{
    \zA{风庭月榭，惜未宴集诗人；帘杏溪桃，或可醉飞吟盏。}
}
{
    \eA{in which they had to struggle for reputation and to fight for gain, but that
        they nevertheless acquired spots with hills and dripping streams,}
}
{
}

\Verse{18}
{
    \zA{孰谓雄才莲社，独许须眉；不教雅会东山，让余脂粉耶？}
}
{
    \eA{and, inviting people to come from far and near, they did all they could to
        detain them, by throwing the linch-pins of their chariots into wells or by
        holding on to their shafts;}
}
{
}

\Verse{19}
{
    \zA{若蒙造雪而来，敢请扫花 以俟。}
}
{
    \eA{and that they invariably joined friendship with two or three of the same
        mind as themselves, with whom they strolled about in these grounds,}
}
{
}

\Verse{20}
{
    \zA{谨启。}
}
{
    \eA{either erecting altars for song,}
}
{
}

\Verse{21}
{
    \zA{宝玉看了，不觉喜的拍手笑道：“倒是三妹妹高雅，我如今就去商议。”}
}
{
    \eA{or establishing societies for scanning poetical works. Their meetings were,
        it is true, prompted, on the spur of the moment, by a sudden fit of good
        cheer, but these have again and again proved,}
}
{
}

\Verse{22}
{
    \zA{一面说， 一面就走。}
}
{
    \eA{during many years, a pleasant topic of conversation.}
}
{
}

\Verse{23}
{
    \zA{翠墨跟在后面。}
}
{
    \eA{I, your cousin, may, I admit,}
}
{
}

\Verse{24}
{
    \zA{刚到了沁芳亭，只见园中后门上值日的婆子手里拿着一个字帖儿走来，见了宝 玉，便迎上去，口内说道：“芸哥儿请安，在后门等着呢。}
}
{
    \eA{be devoid of talent, yet I have been fortunate enough to enjoy your company
        amidst streams and rockeries, and to furthermore admire the elegant verses
        composed by Hsüeh Pao-ch'ai and Lin Tai-yü. When we were in the breezy hall
        and the moonlit pavilion, what a pity we never talked about poets! But near
        the almond tree with the sign and the peach tree by the stream, we may
        perhaps, when under the fumes of wine,}
}
{
}

\Verse{25}
{
    \zA{这是叫我送来的。”}
}
{
    \eA{be able to fling round the cups, used for humming verses!}
}
{
}

\Verse{26}
{
    \zA{宝 玉打开看时，写道： 不肖男芸恭请父亲大人万福金安：男思自蒙天恩，认于膝下，日夜思一孝顺， 竟无可孝顺之处。}
}
{
    \eA{Who is it who opines that societies with any claim to excellent abilities
        can only be formed by men? May it not be that the pleasant meetings on the
        Tung Shan might yield in merit to those, such as ourselves, of the weaker
        sex? Should you not think it too much to walk on the snow,}
}
{
}

\Verse{27}
{
    \zA{前因买办花草，上托大人洪福，竟认得许多花儿匠，并认得许多 名园。}
}
{
    \eA{I shall make bold to ask you round, and sweep the way clean of flowers and
        wait for you. Respectfully written." The perusal of this note filled Pao-yü
        unawares with exultation.}
}
{
}

\Verse{28}
{
    \zA{前因忽见有白海棠一种，不可多得，故变尽方法，只弄得两盆。}
}
{
    \eA{Clapping his hands; "My third cousin," he laughed, "is the one eminently
        polished; I'll go at once to-day and talk matters over with her."}
}
{
}

\Verse{29}
{
    \zA{大人若视男 是亲男一般，便留下赏玩。}
}
{
    \eA{As he spoke, he started immediately, followed by Ts'ui Mo. As soon as they
        reached the Hsin Fang pavilion,}
}
{
}

\Verse{30}
{
    \zA{因天气暑热，恐园中姑娘们妨碍不便，故不敢面见。}
}
{
    \eA{they espied the matron, on duty that day at the back door of the garden,
        advancing towards them with a note in her hand. The moment she perceived
        Pao-yü she forthwith came up to meet him.}
}
{
}

\Verse{31}
{
    \zA{谨 奉书恭启，并叩台安。}
}
{
    \eA{"Mr. Yün," she said, "presents his compliments to you. He is waiting for you
        at the back gate.}
}
{
}

\Verse{32}
{
    \zA{男芸跪书。}
}
{
    \eA{This is a note he bade me bring you."}
}
{
}

\Verse{33}
{
    \zA{宝玉看了，笑问道：“他独来了，还有什么人？”}
}
{
    \eA{Upon opening the note, Pao-yü found it to read as follows: "An unfilial son,
        Yün, reverently inquires about his worthy father's boundless happiness and
        precious health.}
}
{
}

\Verse{34}
{
    \zA{婆子道：“还有两盆花儿。”}
}
{
    \eA{Remembering the honour conferred upon me by your recognising me, in your
        heavenly bounty,}
}
{
}

\Verse{35}
{
    \zA{宝 玉道：“你出去说：我知道了，难为他想着。}
}
{
    \eA{as your son, I tried both day as well as night to do something in evidence
        of my pious obedience, but no opportunity could I find to perform anything
        filial.}
}
{
}

\Verse{36}
{
    \zA{你就把花儿送到我屋里去就是了。”}
}
{
    \eA{When I had, some time back, to purchase flowers and plants, I succeeded,
        thanks to your vast influence,}
}
{
}

\Verse{37}
{
    \zA{一面说，一面同翠墨往秋爽斋来，只见宝钗、黛玉、迎春、惜春已都在那里了。}
}
{
    \eA{venerable senior, in finally making friends with several gardeners and in
        seeing a good number of gardens. As the other day I unexpectedly came across
        a white begonia, of a rare species, I exhausted every possible means to get
        some and managed to obtain just two pots.}
}
{
}

\Verse{38}
{
    \zA{众人见他进来，都大笑说：“又来了一个。”}
}
{
    \eA{If you, worthy senior, regard your son as your own very son, do keep them to
        feast your eyes upon! But with this hot weather to-day,}
}
{
}

\Verse{39}
{
    \zA{探春笑道：“我不算俗，偶然起了个 念头，写了几个帖儿试一试，谁知一招皆到。”}
}
{
    \eA{the young ladies in the garden will, I fear, not be at their ease. I do not
        consequently presume to come and see you in person, so I present you this
        letter, written with due respect,}
}
{
}

\Verse{40}
{
    \zA{宝玉笑道：“可惜迟了!早该起个 社的。”}
}
{
    \eA{while knocking my head before your table. Your son, Yün, on his knees, lays
        this epistle at your feet. A joke!" After reading this note,}
}
{
}

\Verse{41}
{
    \zA{黛玉说道：“此时还不算迟，也没什么可惜；但只你们只管起社，可别算 我，我是不敢的。”}
}
{
    \eA{Pao-yü laughed. "Has he come alone?" he asked. "Or has he any one else with
        him?" "He's got two flower pots as well," rejoined the matron. "You go and
        tell him," Pao-yü urged, "that I've informed myself of the contents of his
        note,}
}
{
}

\Verse{42}
{
    \zA{迎春笑道：“你不敢，谁还敢呢？”}
}
{
    \eA{and that there are few who think of me as he does! If you also take the
        flowers and,}
}
{
}

\Verse{43}
{
    \zA{宝玉道：“这是一件正经 大事，大家鼓舞起来，别你谦我让的。}
}
{
    \eA{put them in my room, it will be all right." So saying, he came with Ts'ui Mo
        into the Ch'iu Shuang study, where he discovered Pao-ch'ai, Tai-yü, Ying
        Ch'un and Hsi Ch'un already assembled.}
}
{
}

\Verse{44}
{
    \zA{各有主意只管说出来，大家评论。}
}
{
    \eA{When they saw him drop in upon them, they all burst out laughing. "Here
        comes still another!"}
}
{
}

\Verse{45}
{
    \zA{宝姐姐也 出个主意，林妹妹也说句话儿。”}
}
{
    \eA{they exclaimed. "I'm not a boor," smiled T'an Ch'un, "so when the idea
        casually crossed my mind, I wrote a few notes to try and see who would come.}
}
{
}

\Verse{46}
{
    \zA{宝钗道：“你忙什么!人还不全呢。”}
}
{
    \eA{But who'd have thought that, as soon as I asked you, you would all come."
        "It's unfortunately late,"}
}
{
}

\Verse{47}
{
    \zA{一语未了， 李纨也来了，进门笑道：“雅的很哪!要起诗社，我自举我掌坛。}
}
{
    \eA{Pao-yü smilingly observed. "We should have started this society long ago."
        "You can't call this late!" Tai-yü interposed, "so why give way to regret!}
}
{
}

\Verse{48}
{
    \zA{前儿春天，我原 有这个意思的，我想了一想，我又不会做诗，瞎闹什么，因而也忘了，就没有说。}
}
{
    \eA{The only thing is, you must form your society, without including me in the
        number; for I daren't be one of you." "If you daren't," Ying Ch'un smiled,
        "who can presume to do so?"}
}
{
}

\Verse{49}
{
    \zA{即是三妹妹高兴，我就帮着你作兴起来。”}
}
{
    \eA{"This is," suggested Pao-yü, "a legitimate and great purpose; and we should
        all exert our energies. You shouldn't be modest,}
}
{
}

\Verse{50}
{
    \zA{黛玉道：“既然定要起诗社，咱们就是诗翁了，先把这些‘姐妹叔嫂’的字样 改了才不俗。”}
}
{
    \eA{and I yielding; but every one of us, who thinks of anything, should freely
        express it for general discussion. So senior cousin Pao-ch'ai do make some
        suggestion; and you junior cousin Lin Tai-yü say something." "What are you
        in this hurry for?"}
}
{
}

\Verse{51}
{
    \zA{李纨道：“极是。}
}
{
    \eA{Pao-ch'ai exclaimed. "We are not all here yet."}
}
{
}

\Verse{52}
{
    \zA{何不起个别号，彼此称呼倒雅？}
}
{
    \eA{This remark was barely concluded, when Li Wan also arrived. As soon as she
        crossed the threshold,}
}
{
}

\Verse{53}
{
    \zA{我是定了‘稻香 老农’，再无人占的。”}
}
{
    \eA{"It's an excellent proposal," she laughingly cried, "this of starting a
        poetical society.}
}
{
}

\Verse{54}
{
    \zA{探春笑道：“我就是‘秋爽居士’罢。”}
}
{
    \eA{I recommend myself as controller. Some time ago in spring, I thought of
        this, 'but,'}
}
{
}

\Verse{55}
{
    \zA{宝玉道：“‘居 士’‘主人’，到底不雅，又累赘。}
}
{
    \eA{I mused, 'I am unable to compose verses, so what's the use of making a mess
        of things?' This is why I dispelled the idea from my mind,}
}
{
}

\Verse{56}
{
    \zA{这里梧桐芭蕉尽有，或指桐蕉起个倒好。”}
}
{
    \eA{and made no mention about it. But since it's your good pleasure, cousin
        Tertia, to start it, I'll help you to set it on foot."}
}
{
}

\Verse{57}
{
    \zA{探 春笑道：“有了，我却爱这芭蕉，就称‘蕉下客’罢。”}
}
{
    \eA{"As you've made up your minds," Tai-yü put in, "to initiate a poetical
        society, every one of us will be poets, so we should, as a first step,}
}
{
}

\Verse{58}
{
    \zA{众人都道别致有趣。}
}
{
    \eA{do away with those various appellations of cousin and uncle and aunt,}
}
{
}

\Verse{59}
{
    \zA{黛玉 笑道：“你们快牵了他来，炖了肉脯子来吃酒。”}
}
{
    \eA{and thus avoid everything that bears a semblance of vulgarity." "First
        rate," exclaimed Li Wan, "and why should we not fix upon some new
        designations by which to address ourselves?}
}
{
}

\Verse{60}
{
    \zA{众人不解，黛玉笑道：“庄子说 的‘蕉叶覆鹿’，他自称‘蕉下客’，可不是一只鹿么？}
}
{
    \eA{This will be a far more refined way! As for my own, I've selected that of
        the 'Old farmer of Tao Hsiang;' so let none of you encroach on it." "I'll
        then call myself the 'resident-scholar of the Ch'iu Shuang,'}
}
{
}

\Verse{61}
{
    \zA{快做了鹿脯来。”}
}
{
    \eA{and have done," T'an Ch'un observed with a smile.}
}
{
}

\Verse{62}
{
    \zA{众人听 了都笑起来。}
}
{
    \eA{"'Resident-scholar or master' is, in fact, not to the point.}
}
{
}

\Verse{63}
{
    \zA{探春因笑道：“你又使巧话来骂人!你别忙，我已替你想了个极当的 美号了。”}
}
{
    \eA{It's clumsy, besides," Pao-yü interposed. "The place here is full of
        dryandra and banana trees, and if one could possibly hit upon some name
        bearing upon the dryandra and banana, it would be preferable."}
}
{
}

\Verse{64}
{
    \zA{又向众人道：“当日娥皇女英洒泪竹上成斑，故今斑竹又名湘妃竹。}
}
{
    \eA{"I've got one," shouted T'an Ch'un smilingly. "I'll style myself 'the guest
        under the banana trees.'" "How uncommon!" they unanimously cried. "It's a
        nice one!" "You had better,"}
}
{
}

\Verse{65}
{
    \zA{如 今他住的是潇湘馆，他又爱哭，将来他那竹子想来也是要变成斑竹的，以后都叫他 做‘潇湘妃子’就完了。”}
}
{
    \eA{laughed Tai-yü, "be quick and drag her away and stew some slices of her
        flesh, for people to eat with their wine." No one grasped her meaning,
        "Ch'uang-tzu," Tai-yü proceeded to explain, smiling, "says: 'The banana
        leaves shelter the deer,' and as she styles herself the guest under the
        banana tree,}
}
{
}

\Verse{66}
{
    \zA{大家听说都拍手叫妙，黛玉低了头也不言语。}
}
{
    \eA{is she not a deer? So be quick and make pieces of dried venison of her." At
        these words, the whole company laughed. "Don't be in a hurry!" T'an Ch'un
        remarked,}
}
{
}

\Verse{67}
{
    \zA{李纨笑道： “我替薛大妹妹也早已想了个好的，也只三个字。”}
}
{
    \eA{as she laughed. "You make use of specious language to abuse people; but I've
        thought of a fine and most apposite name for you!" Whereupon addressing
        herself to the party,}
}
{
}

\Verse{68}
{
    \zA{众人忙问是什么，李纨道：“我 是封他为‘蘅芜君’，不知你们以为如何？”}
}
{
    \eA{"In days gone by," she added, "an imperial concubine, Nü Ying, sprinkled her
        tears on the bamboo, and they became spots, so from olden times to the
        present spotted bamboos have been known as the 'Hsiang imperial concubine
        bamboo.'}
}
{
}

\Verse{69}
{
    \zA{探春道：“这个封号极好。”}
}
{
    \eA{Now she lives in the Hsiao Hsiang lodge, and has a weakness too for tears,}
}
{
}

\Verse{70}
{
    \zA{宝玉道：“我呢？}
}
{
    \eA{so the bamboos over there will by and bye,}
}
{
}

\Verse{71}
{
    \zA{你们也替我想一个。”}
}
{
    \eA{I presume, likewise become transformed into speckled bamboos;}
}
{
}

\Verse{72}
{
    \zA{宝钗笑道：“你的号早有了：‘无事 忙。’}
}
{
    \eA{every one therefore must henceforward call her the 'Hsiao Hsiang imperial
        concubine' and finish with it."}
}
{
}

\Verse{73}
{
    \zA{三字恰当得很！”}
}
{
    \eA{After listening to her, they one and all clapped their hands,}
}
{
}

\Verse{74}
{
    \zA{李纨道：“你还是你的旧号‘绛洞花主’就是了。”}
}
{
    \eA{and cried out: "Capital!" Lin Tai-yü however drooped her head and did not so
        much as utter a single word. "I've also," Li Wan smiled,}
}
{
}

\Verse{75}
{
    \zA{宝玉 笑道：“小时候干的营生，还提他做什么。”}
}
{
    \eA{"devised a suitable name for senior cousin, Hsüeh Pao-chai. It too is one of
        three characters." "What's it?" eagerly inquired the party.}
}
{
}

\Verse{76}
{
    \zA{宝钗道：“还是我送你个号罢，有最 俗的一个号，却于你最当：天下难得的是富贵，又难得的是闲散，这两样再不能兼， 不想你兼有了，就叫你‘富贵闲人’也罢了。”}
}
{
    \eA{"I'll raise her to the rank of 'Princess of Heng Wu,'" Li Wan rejoined. "I
        wonder what you all think about this." "This title of honour," T'an Ch'un
        observed, "is most apposite." "What about mine?" Pao-yü asked. "You should
        try and think of one for me also!"}
}
{
}

\Verse{77}
{
    \zA{宝玉笑道：“当不起，当不起!倒 是随你们混叫去罢。”}
}
{
    \eA{"Your style has long ago been decided upon," Pao-ch'ai smiled. "It consists
        of three words: 'fussing for nothing!' It's most pat!"}
}
{
}

\Verse{78}
{
    \zA{黛玉道：“混叫如何使得!你既住怡红院，索性叫‘怡红公 子’不好？”}
}
{
    \eA{"You should, after all, retain your old name of 'master of the flowers in
        the purple cave,'" Li Wan suggested. "That will do very well." "Those were
        some of the doings of my youth; why rake them up again?"}
}
{
}

\Verse{79}
{
    \zA{众人道：“也好。”}
}
{
    \eA{Pao-yü laughed. "Your styles are very many,"}
}
{
}

\Verse{80}
{
    \zA{李纨道：“二姑娘、四姑娘起个什么？”}
}
{
    \eA{T'an Ch'un observed, "and what do you want to choose another for? All you've
        got to do is to make suitable reply when we call you whatever takes our
        fancy." "I must however give you a name,"}
}
{
}

\Verse{81}
{
    \zA{迎春道： “我们又不大会诗，白起个号做什么！”}
}
{
    \eA{Pao-ch'ai remarked. "There's a very vulgar name, but it's just the very
        thing for you. What is difficult to obtain in the world are riches and
        honours;}
}
{
}

\Verse{82}
{
    \zA{探春道：“虽如此，也起个才是。”}
}
{
    \eA{what is not easy to combine with them is leisure. These two blessings cannot
        be enjoyed together,}
}
{
}

\Verse{83}
{
    \zA{宝钗 道：“他住的是紫菱洲，就叫他‘菱洲’；四丫头住藕香榭，就叫他‘藕榭’就完 了。”}
}
{
    \eA{but, as it happens, you hold one along with the other, so that we might as
        well dub you the 'rich and honourable idler.'" "It won't do; it isn't
        suitable," Pao-yü laughed. "It's better that you should call me, at random,
        whatever you like."}
}
{
}

\Verse{84}
{
    \zA{李纨道：“就是这样好。}
}
{
    \eA{"What names are to be chosen for Miss Secunda and Miss Quarta?"}
}
{
}

\Verse{85}
{
    \zA{但序齿我大，你们都要依我的主意，管教说了大家合 意。}
}
{
    \eA{Li Wan inquired. "We also don't excel in versifying; what's the use
        consequently of giving us names, all for no avail?" Ying Ch'un said. "In
        spite of this," argued T'an Ch'un,}
}
{
}

\Verse{86}
{
    \zA{我们七个人起社，我和二姑娘四姑娘都不会做诗，须得让出我们三个人去。}
}
{
    \eA{"it would be well to likewise find something for you!" "She lives in the Tzu
        Ling Chou, (purple caltrop Isle), so let us call her 'Ling Chou,'" Pao-ch'ai
        suggested. "As for that girl Quarta, she lives in the On Hsiang Hsieh,}
}
{
}

\Verse{87}
{
    \zA{我 们三个人各分一件事。”}
}
{
    \eA{(lotus fragrance pavilion); she should thus be called On Hsieh and have
        done!"}
}
{
}

\Verse{88}
{
    \zA{探春笑道：“已有了号，还只管这样称呼，不如不有了。}
}
{
    \eA{"These will do very well!" Li Wan cried. "But as far as age goes, I am the
        senior, and you should all defer to my wishes; but I feel certain that when
        I've told you what they are,}
}
{
}

\Verse{89}
{
    \zA{以后错了，也要立个罚约才好。”}
}
{
    \eA{you will unanimously agree to them. We are seven here to form the society,
        but neither I, nor Miss Secunda,}
}
{
}

\Verse{90}
{
    \zA{李纨道：“立定了社，再定罚约。}
}
{
    \eA{nor Miss Quarta can write verses; so if you will exclude us three, we'll
        each share some special duties."}
}
{
}

\Verse{91}
{
    \zA{我那里地方儿 大，竟在我那里作社，我虽不能做诗，这些诗人竟不厌俗，容我做个东道主人，我 自然也清雅起来了；还要推我做社长。}
}
{
    \eA{"Their names have already been chosen," T'an Ch'un smilingly demurred; "and
        do you still keep on addressing them like this? Well, in that case, won't it
        be as well for them to have no names? But we must also decide upon some
        scale of fines, for future guidance, in the event of any mistakes."}
}
{
}

\Verse{92}
{
    \zA{我一个社长自然不够，必要再请两位副社长， 就请菱洲藕榭二位学究来，一位出题限韵，一位誊录监场。}
}
{
    \eA{"There will be ample time to fix upon a scale of fines after the society has
        been definitely established." Li Wan replied. "There's plenty of room over
        in my place so let's hold our meetings there. I'm not, it is true, a good
        hand at verses, but if you poets won't treat me disdainfully as a rustic
        boor,}
}
{
}

\Verse{93}
{
    \zA{亦不可拘定了我们三个 不做，若遇见容易些的题目韵脚，我们也随便做一首，你们四个却是要限定的。}
}
{
    \eA{and if you will allow me to play the hostess, I may certainly also gradually
        become more and more refined. As for conceding to me the presidentship of
        the society, it won't be enough, of course, for me alone to preside; it will
        be necessary to invite two others to serve as vice-presidents;}
}
{
}

\Verse{94}
{
    \zA{是 这么着就起，若不依我，我也不敢附骥了。”}
}
{
    \eA{you might then enlist Ling Chou and Ou Hsieh, both of whom are cultured
        persons. The one to choose the themes and assign the metre,}
}
{
}

\Verse{95}
{
    \zA{迎春惜春本性懒于诗词，又有薛林在 前，听了这话，深合己意，二人皆说：“是极。”}
}
{
    \eA{the other to act as copyist and supervisor. We three cannot, however,
        definitely say that we won't write verses, for, if we come across any
        comparatively easy subject and metre,}
}
{
}

\Verse{96}
{
    \zA{探春等也知此意，见他二人悦服， 也不好相强，只得依了。}
}
{
    \eA{we too will indite a stanza if we feel so disposed. But you four will
        positively have to do so. If you agree to this, well, we can proceed with
        the society;}
}
{
}

\Verse{97}
{
    \zA{因笑道：“这话罢了。}
}
{
    \eA{but, if you don't fall in with my wishes, I can't presume to join you."}
}
{
}

\Verse{98}
{
    \zA{只是自想好笑，好好儿的我起了个 主意，反叫你们三个管起我来了。”}
}
{
    \eA{Ying Ch'un and Hsi Ch'un had a natural aversion for verses. What is more,
        Hsüeh Pao-ch'ai and Lin Tai-yü were present. As soon therefore as they heard
        these proposals, which harmonised so thoroughly with their own views,}
}
{
}

\Verse{99}
{
    \zA{宝玉道：“既这样，咱们就往稻香村去。”}
}
{
    \eA{they both, with one voice, approved them as excellent. T'an Ch'un and the
        others were likewise well aware of their object,}
}
{
}

\Verse{100}
{
    \zA{李纨道：“都是你忙。}
}
{
    \eA{but they could not, when they saw with what willingness they accepted the
        charge insist,}
}
{
}

\Verse{101}
{
    \zA{今日不过商 议了，等我再请。”}
}
{
    \eA{with any propriety, upon their writing verses, and they felt obliged to say
        yes.}
}
{
}

\Verse{102}
{
    \zA{宝钗道：“也要议定几日一会才好。”}
}
{
    \eA{"Your proposals," she consequently said, "may be right enough; but in my
        views they are ridiculous.}
}
{
}

\Verse{103}
{
    \zA{探春道：“若只管会多 了，又没趣儿了。}
}
{
    \eA{For here I've had the trouble of initiating this idea of a society, and,
        instead of my having anything to say in the matter,}
}
{
}

\Verse{104}
{
    \zA{一月之中，只可两三次。”}
}
{
    \eA{I've been the means of making you three come and exercise control over me."}
}
{
}

\Verse{105}
{
    \zA{宝钗说道：“一月只要两次就够了。}
}
{
    \eA{"Well then," Pao-yü suggested, "let's go to the Tao Hsiang village." "You're
        always in a hurry!"}
}
{
}

\Verse{106}
{
    \zA{拟定日期，风雨无阻。}
}
{
    \eA{Li Wan remarked. "We're here to-day to simply deliberate. So wait until I've
        sent for you again."}
}
{
}

\Verse{107}
{
    \zA{除这两日外，倘有高兴的，他情愿加一社，或请到他那里去， 或附就了来，也使得。}
}
{
    \eA{"It would be well," Pao-ch'ai interposed, "that we should also decide every
        how many days we are to meet." "If we meet too often," argued T'an Ch'un,
        "there won't be fun in it. We should simply come together two or three times
        in a month."}
}
{
}

\Verse{108}
{
    \zA{岂不活泼有趣？”}
}
{
    \eA{"It will be ample if we meet twice or thrice a month,"}
}
{
}

\Verse{109}
{
    \zA{众人都道：“这个主意更好。”}
}
{
    \eA{Pao-ch'ai added. "But when the dates have been settled neither wind nor rain
        should prevent us.}
}
{
}

\Verse{110}
{
    \zA{探春道： “这原是我起的意，我须得先做个东道，方不负我这番高兴。”}
}
{
    \eA{Exclusive, however, of these two days, any one in high spirits and disposed
        to have an extra meeting can either ask us to go over to her place, or you
        can all come to us; either will do well enough!}
}
{
}

\Verse{111}
{
    \zA{李纨道：“既这样 说，明日你就先开一社不好吗？”}
}
{
    \eA{But won't it be more pleasant if no hard-and-fast dates were laid down?"
        "This suggestion is excellent," they all exclaimed. "This idea was primarily
        originated by me,"}
}
{
}

\Verse{112}
{
    \zA{探春道：“明日不如今日，就是此刻好。}
}
{
    \eA{T'an Ch'un observed, "and I should be the first to play the hostess, so that
        these good spirits of mine shouldn't all go for nothing."}
}
{
}

\Verse{113}
{
    \zA{你就出 题，菱洲限韵，藕榭监场。”}
}
{
    \eA{"Well, after this remark," Li Wan proceeded, "what do you say to your being
        the first to convene a meeting to-morrow?"}
}
{
}

\Verse{114}
{
    \zA{迎春道：“依我说，也不必随一人出题限韵，竟是拈 阄儿公道。”}
}
{
    \eA{"To-morrow," T'an Ch'un demurred, "is not as good as to-day; the best thing
        is to have it at once! You'd better therefore choose the subjects, while
        Ling Chou can fix the metre,}
}
{
}

\Verse{115}
{
    \zA{李纨道：“方才我来时，看见他们抬进两盆白海棠来，倒很好，你们 何不就咏起他来呢？”}
}
{
    \eA{and Ou Hsieh act as supervisor." "According to my ideas," Ying Ch'un chimed
        in, "we shouldn't yield to the wishes of any single person in the choice of
        themes and the settlement of the rhythm. What would really be fair and right
        would be to draw lots."}
}
{
}

\Verse{116}
{
    \zA{迎春道：“都还未赏，先倒做诗？”}
}
{
    \eA{"When I came just now," Li Wan pursued, "I noticed them bring in two pots of
        white begonias,}
}
{
}

\Verse{117}
{
    \zA{宝钗道：“不过是白海 棠，又何必定要见了才做。}
}
{
    \eA{which were simply beautiful; and why should you not write some verses on
        them?" "Can we write verses," Ying Ch'un retorted,}
}
{
}

\Verse{118}
{
    \zA{古人的诗赋也不过都是寄兴寓情，要等见了做，如今也 没这些诗了。”}
}
{
    \eA{"before we have as yet seen anything of the flowers?" "They're purely and
        simply white begonias," Pao-chai answered, "and is there again any need to
        see them before you put together your verses?}
}
{
}

\Verse{119}
{
    \zA{迎春道：“这么着，我就限韵了。”}
}
{
    \eA{Men of old merely indited poetical compositions to express their good cheer
        and conceal their sentiments;}
}
{
}

\Verse{120}
{
    \zA{说着，走到书架前，抽出一本 诗来随手一揭。}
}
{
    \eA{had they waited to write on things they had seen, why, the whole number of
        their works would not be in existence at present!"}
}
{
}

\Verse{121}
{
    \zA{这首诗竟是一首七言律，递与众人看了，都该做七言律。}
}
{
    \eA{"In that case," Ying Ch'un said, "let me fix the metre." With these words,
        she walked up to the book-case, and, extracting a volume, she opened it, at
        random,}
}
{
}

\Verse{122}
{
    \zA{迎春掩了 诗，又向一个小丫头道：“你随口说个字来。”}
}
{
    \eA{at some verses which turned out to be a heptameter stanza. Then handing it
        round for general perusal, everybody had to compose lines with seven words
        in each.}
}
{
}

\Verse{123}
{
    \zA{那丫头正倚门站着，便说了个“门” 字，迎春笑道：“就是‘门’字韵，‘十三元’了。}
}
{
    \eA{Ying Ch'un next closed the book of verses and addressed herself to a young
        waiting-maid. "Just utter," she bade her, "the first character that comes to
        your mouth." The waiting-maid was standing,}
}
{
}

\Verse{124}
{
    \zA{起头一个韵定要‘门’字。”}
}
{
    \eA{leaning against the door, so readily she suggested the word "door."}
}
{
}

\Verse{125}
{
    \zA{说着又要了韵牌匣子过来，抽出“十三元”一屉，又命那丫头随手拿四块。}
}
{
    \eA{"The rhyme then will be the word 'door,'" Ying Ch'un smiled, "under the
        thirteenth character 'Yuan.' The final word of the first line is therefore
        'door'." Saying this, she asked for the box with the rhyme slips,}
}
{
}

\Verse{126}
{
    \zA{那丫头 便拿了“盆”“魂”“痕”“昏”四块来。}
}
{
    \eA{and, pulling out the thirteenth drawer with the character "Yuan," she
        directed a young waiting-maid to take four words as they came under her
        hand.}
}
{
}

\Verse{127}
{
    \zA{宝玉道：“这‘盆’‘门’两个字不大 好做呢！”}
}
{
    \eA{The waiting-maid complied with her directions, and picked out four slips, on
        which were written "p'en, hun, hen and hun," pot, spirit,}
}
{
}

\Verse{128}
{
    \zA{侍书一样预备下四分纸笔，便都悄然各自思索起来。}
}
{
    \eA{traces and dusk. "The two characters pot and door," observed Pao-yü, "are
        not very easy to rhyme with." But Shih Shu then got ready four lots of paper
        and pens,}
}
{
}

\Verse{129}
{
    \zA{独黛玉或抚弄梧桐，或看 秋色，或又和丫鬟们嘲笑。}
}
{
    \eA{share and share alike, and one and all quietly set to work, racking their
        brains to perform their task, with the exception of Tai-yü,}
}
{
}

\Verse{130}
{
    \zA{迎春又命丫鬟点了一枝梦甜香。}
}
{
    \eA{who either kept on rubbing the dryandra flowers, or looking at the autumnal
        weather,}
}
{
}

\Verse{131}
{
    \zA{原来这梦甜香只有三寸 来长，有灯草粗细，以其易烬，故以此为限，如香烬未成便要受罚。}
}
{
    \eA{or bandying jokes as well with the servant-girls; while Ying Ch'un ordered a
        waiting-maid to light a "dream-sweet" incense stick. This "dream-sweet"
        stick was, it must be explained, made only about three inches long and about
        the thickness of a lamp-wick,}
}
{
}

\Verse{132}
{
    \zA{一时探春便先 有了，自己提笔写出，又改抹了一回，递与迎春。}
}
{
    \eA{in order to easily burn down. Setting therefore her choice upon one of these
        as a limit of time, any one who failed to accomplish the allotted task, by
        the time the stick was consumed,}
}
{
}

\Verse{133}
{
    \zA{因问宝钗：“蘅芜君，你可有了？”}
}
{
    \eA{had to pay a penalty. Presently, T'an Ch'un was the first to think of some
        verses,}
}
{
}

\Verse{134}
{
    \zA{宝钗道：“有却有了，只是不好。”}
}
{
    \eA{and, taking up her pen, she wrote them down; and, after submitting them to
        several alterations,}
}
{
}

\Verse{135}
{
    \zA{宝玉背着手在回廊上踱来踱去，因向黛玉说道： “你听他们都有了。”}
}
{
    \eA{she handed them up to Ying Ch'un. "Princess of Heng Wu," she then inquired
        of Pao-ch'ai, "have you finished?" "As for finishing, I have finished,"
        Pao-ch'ai rejoined; "but they're worth nothing."}
}
{
}

\Verse{136}
{
    \zA{黛玉道：“你别管我。”}
}
{
    \eA{Pao-yü paced up and down the verandah with his hands behind his back.}
}
{
}

\Verse{137}
{
    \zA{宝玉又见宝钗已誊写出来，因说道： “了不得，香只剩下一寸了!我才有了四句。”}
}
{
    \eA{"Have you heard?" he thereupon said to Tai-yü, "they've all done!" "Don't
        concern yourself about me!" Tai-yü returned for answer. Pao-yü also
        perceived that Pao-ch'ai had already copied hers out.}
}
{
}

\Verse{138}
{
    \zA{又向黛玉道：“香要完了，只管蹲 在那潮地下做什么？”}
}
{
    \eA{"Dreadful!" he exclaimed. "There only remains an inch of the stick and I've
        only just composed four lines. The incense stick is nearly burnt out,"}
}
{
}

\Verse{139}
{
    \zA{黛玉也不理。}
}
{
    \eA{he continued, speaking to Tai-yü,}
}
{
}

\Verse{140}
{
    \zA{宝玉道：“我可顾不得你了，管他好歹，写出 来罢。”}
}
{
    \eA{"and what do you keep squatting on that damp ground like that for?" But
        Tai-yü did not again worry her mind about what he said. "Well," Pao-yü
        added, "I can't be looking after you! Whether good or bad, I'll write mine
        out too and have done."}
}
{
}

\Verse{141}
{
    \zA{说着，走到案前写了。}
}
{
    \eA{As he spoke, he likewise drew up to the table and began putting his lines
        down.}
}
{
}

\Verse{142}
{
    \zA{李纨道：“我们要看诗了。}
}
{
    \eA{"We'll now peruse the verses," Li Wan interposed, "and if by the time we've
        done,}
}
{
}

\Verse{143}
{
    \zA{若看完了还不交卷，是必罚的。”}
}
{
    \eA{you haven't as yet handed up your papers, you'll have to be fined." "Old
        farmer of Tao Hsiang,"}
}
{
}

\Verse{144}
{
    \zA{宝玉道：“稻香 老农虽不善作，却善看，又最公道，你的评阅，我们是都服的。”}
}
{
    \eA{Pao-yü remarked, "you're not, it is true, a good hand at writing verses, but
        you can read well, and, what's more, you're the fairest of the lot; so you'd
        better adjudge the good and bad,}
}
{
}

\Verse{145}
{
    \zA{众人点头。}
}
{
    \eA{and we'll submit to your judgment."}
}
{
}

\Verse{146}
{
    \zA{于是 先看探春的稿上写道： 咏白海棠 斜阳寒草带重门，苔翠盈铺雨后盆。}
}
{
    \eA{"Of course!" responded the party with one voice. In due course, therefore,
        she first read T'an Ch'un's draft. It ran as follows:— Verses on the
        Begonia. What time the sun's rays slant,}
}
{
}

\Verse{147}
{
    \zA{玉是精神难比洁，雪为肌骨易销魂。}
}
{
    \eA{and the grass waxeth cold, close the double doors. After a shower of rain,
        green moss plenteously covers the whole pot.}
}
{
}

\Verse{148}
{
    \zA{芳心一点娇无力，倩影三更月有痕。}
}
{
    \eA{Beauteous is jade, but yet with thee in purity it cannot ever vie. Thy
        frame, spotless as snow,}
}
{
}

\Verse{149}
{
    \zA{莫道缟仙能羽化，多情伴我咏黄昏。}
}
{
    \eA{from admiration easy robs me of my wits Thy fragrant core is like unto a
        dot, so full of grace, so delicate!}
}
{
}

\Verse{150}
{
    \zA{大家看了，称赏一回，又看宝钗的道： 珍重芳姿昼掩门，自携手瓮灌苔盆。}
}
{
    \eA{When the moon reacheth the third watch, thy comely shade begins to show
        itself. Do not tell me that a chaste fairy like thee can take wings and pass
        away. How lovely are thy charms, when in thy company at dusk I sing my lay!}
}
{
}

\Verse{151}
{
    \zA{胭脂洗出秋阶影，冰雪招来露砌魂。}
}
{
    \eA{After she had read them aloud, one and all sang their praise for a time. She
        then took up Pao-ch'ai's,}
}
{
}

\Verse{152}
{
    \zA{淡极始知花更艳，愁多焉得玉无痕？}
}
{
    \eA{which consisted of: If thou would'st careful tend those fragrant lovely
        flowers, close of a day the doors,}
}
{
}

\Verse{153}
{
    \zA{欲偿白帝宜清洁，不语婷婷日又昏。}
}
{
    \eA{And with thine own hands take the can and sprinkle water o'er the mossy
        pots. Red, as if with cosmetic washed,}
}
{
}

\Verse{154}
{
    \zA{李纨笑道：“到底是蘅芜君！”}
}
{
    \eA{are the shadows in autumn on the steps. Their crystal snowy bloom invites
        the dew on their spirits to heap itself.}
}
{
}

\Verse{155}
{
    \zA{说着，又看宝玉的道： 秋容浅淡映重门，七节攒成雪满盆。}
}
{
    \eA{Their extreme whiteness mostly shows that they're more comely than all other
        flowers. When much they grieve, how can their jade-like form lack the traces
        of tears?}
}
{
}

\Verse{156}
{
    \zA{出浴太真冰作影，捧心西子玉为魂。}
}
{
    \eA{Would'st thou the god of those white flowers repay? then purity need'st thou
        observe. In silence plunges their fine bloom,}
}
{
}

\Verse{157}
{
    \zA{晓风不散愁千点，宿雨还添泪一痕。}
}
{
    \eA{now that once more day yields to dusk. "After all," observed Li Wan, "it's
        the Princess of Heng Wu,}
}
{
}

\Verse{158}
{
    \zA{独倚画栏如有意，清砧怨笛送黄昏。}
}
{
    \eA{who expresses herself to the point." Next they bestowed their attention on
        the following lines,}
}
{
}

\Verse{159}
{
    \zA{大家看了，宝玉说探春的好。}
}
{
    \eA{composed by Pao-yü:— Thy form in autumn faint reflects against the double
        doors.}
}
{
}

\Verse{160}
{
    \zA{李纨终要推宝钗：“这诗有身分。”}
}
{
    \eA{So heaps the snow in the seventh feast that it filleth thy pots. Thy shade
        is spotless as Tai Chen,}
}
{
}

\Verse{161}
{
    \zA{因又催黛玉。}
}
{
    \eA{when from her bath she hails.}
}
{
}

\Verse{162}
{
    \zA{黛玉道：“你们都有了？”}
}
{
    \eA{Like Hsi Tzu's, whose hand ever pressed her heart, jade-like thy soul. When
        the morn-ushering breeze falls not,}
}
{
}

\Verse{163}
{
    \zA{说着，提笔一挥而就，掷与众人。}
}
{
    \eA{thy thousand blossoms grieve. To all thy tears the evening shower addeth
        another trace.}
}
{
}

\Verse{164}
{
    \zA{李纨等看他写的道： 半卷湘帘半掩门，碾冰为土玉为盆。}
}
{
    \eA{Alone thou lean'st against the coloured rails as if with sense imbued. As
        heavy-hearted as the fond wife, beating clothes, or her that sadly listens
        to the flute, thou mark'st the fall of dusk.}
}
{
}

\Verse{165}
{
    \zA{看了这句，宝玉先喝起彩来，说：“从何处想来！”}
}
{
    \eA{When they had perused his verses, Pao-yü opined that T'an Ch'un's carried
        the palm. Li Wan was, however, inclined to concede to the stanza,}
}
{
}

\Verse{166}
{
    \zA{又看下面道： 偷来梨蕊三分白，借得梅花一缕魂。}
}
{
    \eA{indited by Pao-ch'ai, the credit of possessing much merit. But she then went
        on to tell Tai-yü to look sharp. "Have you all done?" Tai-yü asked.}
}
{
}

\Verse{167}
{
    \zA{众人看了，也都不禁叫好，说：“果然比别人又是一样心肠。”}
}
{
    \eA{So saying, she picked up a pen and completing her task, with a few dashes,
        she threw it to them to look over. On perusal, Li Wan and her companions
        found her verses to run in this strain:— Half rolled the speckled portiere
        hangs,}
}
{
}

\Verse{168}
{
    \zA{又看下面道： 月窟仙人缝缟袂，秋闺怨女拭啼痕。}
}
{
    \eA{half closed the door. Thy mould like broken ice it looks, jade-like thy pot.
        This couplet over, Pao-yü took the initiative and shouted:}
}
{
}

\Verse{169}
{
    \zA{娇羞默默同谁诉？}
}
{
    \eA{"Capital." But he had just had time to inquire where she had recalled them
        to mind from,}
}
{
}

\Verse{170}
{
    \zA{倦倚西风夜已昏。}
}
{
    \eA{when they turned their mind to the succeeding lines:}
}
{
}

\Verse{171}
{
    \zA{众人看了，都道：“是这首为上。”}
}
{
    \eA{Three points of whiteness from the pear petals thou steal'st; And from the
        plum bloom its spirit thou borrowest.}
}
{
}

\Verse{172}
{
    \zA{李纨道：“若论风流别致，自是这首；若 论含蓄浑厚，终让蘅稿。”}
}
{
    \eA{"Splendid!" every one (who heard) them conned over, felt impelled to cry.
        "It is a positive fact," they said, "that her imagination is, compared with
        that of others,}
}
{
}

\Verse{173}
{
    \zA{探春道：“这评的有理。}
}
{
    \eA{quite unique." But the rest of the composition was next considered.}
}
{
}

\Verse{174}
{
    \zA{潇湘妃子当居第二。”}
}
{
    \eA{Its text was: The fairy in Selene's cavity donneth a plain attire.}
}
{
}

\Verse{175}
{
    \zA{李纨道： “怡红公子是压尾，你服不服？”}
}
{
    \eA{The maiden, plunged in autumn grief, dries in her room the prints of tears.
        Winsome she blushes, in silence she's plunged,}
}
{
}

\Verse{176}
{
    \zA{宝玉道：“我的那首原不好，这评的最公。”}
}
{
    \eA{with none a word she breathes; But wearily she leans against the eastern
        breeze, though dusk has long since fall'n.}
}
{
}

\Verse{177}
{
    \zA{又 笑道：“只是蘅潇二首，还要斟酌。”}
}
{
    \eA{"This stanza ranks above all!" they unanimously remarked, after it had been
        read for their benefit.}
}
{
}

\Verse{178}
{
    \zA{李纨道：“原是依我评论，不与你们相干， 再有多说者必罚。”}
}
{
    \eA{"As regards beauty of thought and originality, this stanza certainly
        deserves credit," Li Wan asserted; "but as regards pregnancy and simplicity
        of language,}
}
{
}

\Verse{179}
{
    \zA{宝玉听说，只得罢了。}
}
{
    \eA{it, after all, yields to that of Heng Wu." "This criticism is right."}
}
{
}

\Verse{180}
{
    \zA{李纨道：“从此后，我定于每月初二、 十六这两日开社，出题限韵都要依我。}
}
{
    \eA{T'an Ch'un put in. "That of the Hsiao Hsiang consort must take second
        place." "Yours, gentleman of I Hung," Li Wan pursued, "is the last of the
        lot. Do you agreeably submit to this verdict?" "My stanza," Pao-yü ventured,}
}
{
}

\Verse{181}
{
    \zA{这其间你们有高兴的，只管另择日子补开， 那怕一个月每天都开社我也不管。}
}
{
    \eA{"isn't really worth a straw. Your criticism is exceedingly fair. But," he
        smilingly added, "the two poems, written by Heng Wu and Hsiao Hsiang, have
        still to be discussed."}
}
{
}

\Verse{182}
{
    \zA{只是到了初二、十六这两日，是必往我那里去。”}
}
{
    \eA{"You should," argued Li Wan, "fall in with my judgment; this is no business
        of any of you, so whoever says anything more will have to pay a penalty."}
}
{
}

\Verse{183}
{
    \zA{宝玉道：“到底要起个社名才是。”}
}
{
    \eA{Pao-yü at this reply found that he had no alternative but to drop the
        subject. "I decide that from henceforward,"}
}
{
}

\Verse{184}
{
    \zA{探春道：“俗了又不好，忒新了刁钻古怪也不 好。}
}
{
    \eA{Li Wan proceeded, "we should hold meetings twice every month, on the second
        and sixteenth. In the selection of themes and the settlement of the rhymes,}
}
{
}

\Verse{185}
{
    \zA{可巧才是海棠诗开端，就叫个‘海棠诗社’罢，虽然俗些，因真有此事，也就 不碍了。”}
}
{
    \eA{you'll all have then to do as I wish. But any person who may, during the
        intervals, feel so disposed, will be at perfect liberty to choose another
        day for an extra meeting. What will I care if there's a meeting every day of
        the moon?}
}
{
}

\Verse{186}
{
    \zA{说毕，大家又商议了一回。}
}
{
    \eA{It will be no concern of mine, so long as when the second and sixteenth
        arrive,}
}
{
}

\Verse{187}
{
    \zA{略用些酒果，方各自散去，也有回家的，也 有往贾母王夫人处去的。}
}
{
    \eA{you do, as you're bound to, and come over to my place." "We should, as is
        but right," Pao-yü suggested, "choose some name or other for our society."}
}
{
}

\Verse{188}
{
    \zA{当下无话。}
}
{
    \eA{"Were an ordinary one chosen,}
}
{
}

\Verse{189}
{
    \zA{且说袭人因见宝玉看了字帖儿，便慌慌张张同翠墨去了，也不知何事；后来又 见后门上婆子送了两盆海棠花来。}
}
{
    \eA{it wouldn't be nice," T'an Ch'un explained, "and anything too new-fangled,
        eccentric or strange won't also be quite the thing! As luck would have it,
        we've just started with the poems on the begonia, so let us call it the
        'Begonia Poetical Society.'}
}
{
}

\Verse{190}
{
    \zA{袭人问那里来的，婆子们便将前番原故说了。}
}
{
    \eA{This title is, it's true, somewhat commonplace; but as it's positively based
        on fact, it shouldn't matter." After this proposal of hers,}
}
{
}

\Verse{191}
{
    \zA{袭 人听说，便命他们摆好，让他们在下房里坐了。}
}
{
    \eA{they held further consultation; and partaking of some slight refreshments,
        each of them eventually retired. Some repaired to their quarters.}
}
{
}

\Verse{192}
{
    \zA{自己走到屋里，称了六钱银子封好， 又拿了三百钱走来，都递给那两个婆子道：“这银子赏那抬花儿的小子们。}
}
{
    \eA{Others went to dowager lady Chia's or Madame Wang's apartments. But we will
        leave them without further comment. When Hsi Jen, for we will now come to
        her, perceived Pao-yü peruse the note and walk off in a great flurry,}
}
{
}

\Verse{193}
{
    \zA{这钱你 们打酒喝罢。”}
}
{
    \eA{along with Ts'ui Mo, she was quite at a loss what to make of it.}
}
{
}

\Verse{194}
{
    \zA{那婆子们站起来，眉开眼笑，千恩万谢的不肯受，见袭人执意不收， 方领了。}
}
{
    \eA{Subsequently, she also saw the matrons, on duty at the back gate, bring two
        pots of begonias. Hsi Jen inquired of them where they came from. The women
        explained to her all about them.}
}
{
}

\Verse{195}
{
    \zA{袭人又道：“后门上外头可有该班的小子们？”}
}
{
    \eA{As soon as Hsi Jen heard their reply, she at once desired them to put the
        flowers in their proper places, and asked them to sit down in the lower
        rooms.}
}
{
}

\Verse{196}
{
    \zA{婆子忙应道：“天天有四 个，原预备里头差使的。}
}
{
    \eA{She then entered the house, and, weighing six mace of silver, she wrapped it
        up properly, and fetching besides three hundred cash,}
}
{
}

\Verse{197}
{
    \zA{姑娘有什么差使？}
}
{
    \eA{she came over and handed both the amounts to the two matrons. "This silver,"
        she said,}
}
{
}

\Verse{198}
{
    \zA{我们吩咐去。”}
}
{
    \eA{"is a present for the boys, who carried the flowers;}
}
{
}

\Verse{199}
{
    \zA{袭人笑道：“我有什么 差使。}
}
{
    \eA{and these cash are for you to buy yourselves a cup of tea with." The women
        rose to their feet in such high glee that their eyebrows dilated and their
        eyes smiled; but, though they waxed eloquent in the expression of their deep
        gratitude, they would not accept the money. It was only after they had
        perceived how obstinate Hsi Jen was in not taking it back that they at last
        volunteered to keep it.}
}
{
}

\Verse{200}
{
    \zA{今儿宝二爷要打发人到小侯爷家给史大姑娘送东西去，可巧你们来了，顺便 出去叫后门上小子们雇辆车来，回来你们就往这里拿钱，不用叫他们往前头混碰 去。”}
}
{
    \eA{"Are there," Hsi Jen then inquired, "any servant-boys on duty outside the
        back gate?" "There are four of them every day," answered one of the matrons.
        "They're put there with the sole idea of attending to any orders that might
        be given them from inside. But, Miss, if you've anything to order them to
        do,}
}
{
}

\Verse{201}
{
    \zA{婆子答应着去了。}
}
{
    \eA{we'll go and deliver your message." "What orders can I have to give them?"}
}
{
}

\Verse{202}
{
    \zA{袭人回至房中，拿碟子盛东西与湘云送去。}
}
{
    \eA{Hsi Jen laughed. "Mr. Pao, our master Secundus, was purposing to send some
        one to-day to the young marquis' house to take something over to Miss Shih.}
}
{
}

\Verse{203}
{
    \zA{却见？}
}
{
    \eA{But you come at an opportune moment so you might,}
}
{
}

\Verse{204}
{
    \zA{子上碟子槽儿空着，因回头 见晴雯、秋纹、麝月等都在一处做针黹，袭人问道：“那个缠丝白玛瑙碟子那里去 了？”}
}
{
    \eA{on your way out, tell the servant-boys at the back gate to hire a carriage;
        and on its return you can come here and get the money. But don't let them
        rush recklessly against people in the front part of the compound!" The
        matrons signified their obedience and took their leave. Hsi Jen retraced her
        steps into the house to fetch a tray in which to place the presents intended
        for Shih Hsiang-yün,}
}
{
}

\Verse{205}
{
    \zA{众人见问，你看我，我看你，都想不起来。}
}
{
    \eA{but she discovered the shelf for trays empty. Upon turning round, however,
        she caught sight of Ch'ing Wen,}
}
{
}

\Verse{206}
{
    \zA{半日晴雯笑道：“给三姑娘送荔 枝去了，还没送来呢。”}
}
{
    \eA{Ch'iu Wen, She Yüeh and the other girls, seated together, busy with their
        needlework. "Where is the white cornelian tray with twisted threads gone
        to?"}
}
{
}

\Verse{207}
{
    \zA{袭人道：“家常送东西的家伙多着呢，巴巴儿的拿这个。”}
}
{
    \eA{Hsi Jen asked. At this question, one looked at the one, and the other stared
        at the other, but none of them could remember anything about it.}
}
{
}

\Verse{208}
{
    \zA{晴雯道：“我也这么说，但只那碟子配上鲜荔枝才好看。}
}
{
    \eA{After a protracted lapse of time, Ch'ing Wen smiled. "It was taken to Miss
        Tertia's with a present of lichees," she rejoined, "and it hasn't as yet
        been returned."}
}
{
}

\Verse{209}
{
    \zA{我送去，三姑娘也见了， 说好看，连碟子放着，就没带来。}
}
{
    \eA{"There are plenty of articles," Hsi Jen remarked, "for sending over things
        on ordinary occasions; and do you deliberately go and carry this off?"}
}
{
}

\Verse{210}
{
    \zA{你再瞧那？}
}
{
    \eA{"Didn't I maintain the same thing?"}
}
{
}

\Verse{211}
{
    \zA{子尽上头的一对联珠瓶还没收来呢。”}
}
{
    \eA{Ch'ing Wen retorted. "But so well did this tray match with the fresh lichees
        it contained, that when I took it over,}
}
{
}

\Verse{212}
{
    \zA{秋纹笑道：“提起这瓶来，我又想起笑话儿来了。}
}
{
    \eA{Miss T'an Ch'un herself noticed the fact. 'How splendid,' she said, and lo,
        putting even the tray by, she never had it brought over.}
}
{
}

\Verse{213}
{
    \zA{我们宝二爷说声孝心一动，也孝 敬到二十分：那日见园里桂花，折了两枝，原是自己要插瓶的，忽然想起来，说： ‘这是自己园里才开的新鲜花儿，不敢自己先玩。’}
}
{
    \eA{But, look! hasn't the pair of beaded vases, which stood on the very top of
        that shelf, been fetched as yet?" "The mention of these vases," Ch'iu Wen
        laughed, "reminds me again of a funny incident. Whenever our Mr. Pao-yü's
        filial piety is aroused,}
}
{
}

\Verse{214}
{
    \zA{巴巴儿的把那对瓶拿下来，亲 自灌水插好了，叫个人拿着，亲自送一瓶进老太太，又进一瓶给太太。}
}
{
    \eA{he shows himself filial over and above the highest degree! The other day, he
        espied the olea flowers in the park, and he plucked two twigs. His original
        idea was to place them in a vase for himself, but a sudden thought struck
        him.}
}
{
}

\Verse{215}
{
    \zA{谁知他孝心 一动，连跟的人都得了福了。}
}
{
    \eA{'These are flowers,' he mused, 'which have newly opened in our garden, so
        how can I presume to be the first to enjoy them?'}
}
{
}

\Verse{216}
{
    \zA{可巧那日是我拿去的，老太太见了喜的无可不可，见 人就说：‘到底是宝玉孝顺我，连一枝花儿也想的到。}
}
{
    \eA{And actually taking down that pair of vases, he filled them with water with
        his own hands, put the flowers in, and, calling a servant to carry them, he
        in person took one of the vases into dowager lady Chia's, and then took the
        other to Madame Wang's.}
}
{
}

\Verse{217}
{
    \zA{别人还只抱怨我疼他！’}
}
{
    \eA{But, as it happens, even his attendants reap some benefit, when once his
        filial feelings are stirred up!}
}
{
}

\Verse{218}
{
    \zA{你 们知道老太太素日不大和我说话，有些不入他老人家的眼；那日竟叫人拿几百钱给 我，说我‘可怜见儿的，生的单弱’。}
}
{
    \eA{As luck would have it, the one who carried the vases over on that day was
        myself. The sight of these flowers so enchanted our venerable lady that
        there was nothing that she wouldn't do. 'Pao-yü,' she said to every one she
        met, 'is the one, after all, who shows me much attention. So much so, that
        he has even thought of bringing me a twig of flowers!}
}
{
}

\Verse{219}
{
    \zA{这可是再想不到的福气。}
}
{
    \eA{And yet, the others bear me a grudge on account of the love that I lavish on
        him!'}
}
{
}

\Verse{220}
{
    \zA{几百钱是小事，难 得这个脸面。}
}
{
    \eA{Our venerable mistress, you all know very well, has never had much to say to
        me.}
}
{
}

\Verse{221}
{
    \zA{及至到了太太那里，太太正和二奶奶赵姨奶奶好些人翻箱子，找太太 当日年轻的颜色衣裳，不知要给那一个；一见了，连衣裳也不找了，且看花儿。}
}
{
    \eA{I have all along not been much of a favourite in the old lady's eyes. But on
        that occasion she verily directed some one to give me several hundreds of
        cash. 'I was to be pitied,' she observed, 'for being born with a weak
        physique!' This was, indeed, an unforeseen piece of good luck! The several
        hundreds of cash are a mere trifle; but what's not easy to get is this sort
        of honour!}
}
{
}

\Verse{222}
{
    \zA{又 有二奶奶在傍边凑趣儿，夸宝二爷又是怎么孝顺，又是怎么知好歹，有的没的说了 两车话。}
}
{
    \eA{After that, we went over into Madame Wang's. Madame Wang was, at the time,
        with our lady Secunda, Mrs. Chao, and a whole lot of people; turning the
        boxes topsy-turvey, trying to find some coloured clothes her ladyship had
        worn long ago in her youth,}
}
{
}

\Verse{223}
{
    \zA{当着众人，太太脸上又增了光，堵了众人的嘴，太太越发喜欢了，现成的 衣裳，就赏了我两件。}
}
{
    \eA{so as to give them to some one or other. Who it was, I don't know. But the
        moment she saw us, she did not even think of searching for any clothes, but
        got lost in admiration for the flowers. Our lady Secunda was also standing
        by,}
}
{
}

\Verse{224}
{
    \zA{衣裳也是小事，年年横竖也得，却不像这个彩头。”}
}
{
    \eA{and she made sport of the matter. She extolled our master Pao, for his
        filial piety and for his knowledge of right and wrong; and what with what
        was true and what wasn't,}
}
{
}

\Verse{225}
{
    \zA{晴雯笑道：“呸!好没见世面的小蹄子!那是把好的给了人，挑剩下的才给你， 你还充有脸呢！”}
}
{
    \eA{she came out with two cart-loads of compliments. These things spoken in the
        presence of the whole company so added to Madame Wang's lustre and sealed
        every one's mouth, that her ladyship was more and more filled with
        gratification, and she gave me two ready-made clothes as a present.}
}
{
}

\Verse{226}
{
    \zA{秋纹道：“凭他给谁剩的，到底是太太的恩典。”}
}
{
    \eA{These too are of no consequence; one way or another, we get some every year;
        but nothing can come up to this sort of lucky chance!"}
}
{
}

\Verse{227}
{
    \zA{晴雯道：“要 是我，我就不要。}
}
{
    \eA{"Psha!" Ch'ing Wen ejaculated with a significant smile, "you are indeed a
        mean thing,}
}
{
}

\Verse{228}
{
    \zA{若是给别人剩的给我也罢了，一样这屋里的人，难道谁又比谁高 贵些？}
}
{
    \eA{who has seen nothing of the world! She gave the good ones to others and the
        refuse to you; and do you still pat on all this side?" "No matter whether
        what she gave me was refuse or not," Ch'iu Wen protested, "it's,}
}
{
}

\Verse{229}
{
    \zA{把好的给他，剩的才给我，我宁可不要，冲撞了太太，我也不受这口气！”}
}
{
    \eA{after all, an act of bounty on the part of her ladyship." "Had it been
        myself," Ch'ing Wen pursued, "I would at once have refused them! It wouldn't
        have mattered if she had given me what had been left by some one else;}
}
{
}

\Verse{230}
{
    \zA{秋纹忙问道：“给这屋里谁的？}
}
{
    \eA{but we all stand on an equal footing in these rooms, and is there any one,
        forsooth,}
}
{
}

\Verse{231}
{
    \zA{我因为前日病了几天，家去了，不知是给谁的，好 姐姐，你告诉我知道。”}
}
{
    \eA{so much the more exalted or honorable than the other as to justify her
        taking what is good and bestowing it upon her and giving me what is left? I
        had rather not take them! I might have had to give offence to Madame Wang,}
}
{
}

\Verse{232}
{
    \zA{晴雯道：“我告诉了你，难道你这会子退还太太去不成？”}
}
{
    \eA{but I wouldn't have put up with such a slight!" "To whom did she give any in
        these rooms?" Ch'iu Wen vehemently inquired. "I was unwell and went home for
        several days,}
}
{
}

\Verse{233}
{
    \zA{秋纹笑道：“胡说!我白听了喜欢喜欢，那怕给这屋里的狗剩下的，我只领太太的 恩典，也不管别的事。”}
}
{
    \eA{so that I am not aware to whom any were given. Dear sister, do tell me who
        it is so that I may know." "Were I to tell you," Ch'ing Wen rejoined, "is it
        likely that you would return them at this hour to Madame Wang?" "What
        nonsense," Ch'iu Wen laughed. "Ever since I've heard about it,}
}
{
}

\Verse{234}
{
    \zA{众人听了都笑道：“骂的巧，可不是给了那西洋花点子哈 巴儿了！”}
}
{
    \eA{I've been delighted and happy. No matter if she even bestowed upon me what
        remained from anything given to a dog in these rooms, I would have been
        thankful for her ladyship's kindness. I wouldn't have worried my mind with
        anything else!"}
}
{
}

\Verse{235}
{
    \zA{袭人笑道：“你们这起烂了嘴的!得空儿就拿我取笑打牙儿，一个个不 知怎么死呢！”}
}
{
    \eA{After listening to her, everybody laughed. "Doesn't she know how to jeer in
        fine style!" they ejaculated unanimously; "for weren't they given to that
        foreign spotted pug dog?" "You lot of filthy-tongued creatures!" Hsi Jen
        laughed, "when you've got nothing to do,}
}
{
}

\Verse{236}
{
    \zA{秋纹笑道：“原来姐姐得了!我实在不知道，我陪个不是罢。”}
}
{
    \eA{you make me the scapegoat to crack your jokes, and poke your fun at! But
        what kind of death will, I wonder, each of you have!" "Was it verily you,
        sister,}
}
{
}

\Verse{237}
{
    \zA{袭 人笑道：“少轻狂罢!你们谁取了碟子来是正经。”}
}
{
    \eA{who got them?" Ch'iu Wen asked with a smile. "I assure you I had no idea
        about it! I tender you my apologies." "You might be a little less
        domineering!"}
}
{
}

\Verse{238}
{
    \zA{麝月道：“那瓶也该得空儿收 来了。}
}
{
    \eA{Hsi Jen remarked smilingly. "The thing now is, who of you will go and fetch
        the tray." "The vases too," Shih Yüeh suggested,}
}
{
}

\Verse{239}
{
    \zA{老太太屋里还罢了，太太屋里人多手杂，别人还可已，那个主儿的一伙子人 见是这屋里的东西，又该使黑心弄坏了才罢。}
}
{
    \eA{"must be got back when there's any time to spare; for there's nothing to say
        about our venerable mistress' quarters, but Madame Wang's apartments teem
        with people and many hands. The rest are all right; but Mrs. Chao and all
        that company will, when they see that the vase hails from these rooms,}
}
{
}

\Verse{240}
{
    \zA{太太又不大管这些，不如早收来是正 经。”}
}
{
    \eA{surely again foster evil designs, and they won't feel happy until they've
        done all they can to spoil it!}
}
{
}

\Verse{241}
{
    \zA{晴雯听说，便放下针线道：“这是等我取去呢。”}
}
{
    \eA{Besides, Madame Wang doesn't trouble herself about such things. So had we
        not as well bring it over a moment sooner?"}
}
{
}

\Verse{242}
{
    \zA{秋纹道：“还是我取去罢， 你取你的碟子去。”}
}
{
    \eA{Hearing this, Ch'ing Wen threw down her needlework. "What you say is
        perfectly right," she assented, "so you'd better let me go and fetch it."}
}
{
}

\Verse{243}
{
    \zA{晴雯道：“我偏取一遭儿。}
}
{
    \eA{"I'll, after all, go for it." Ch'iu Wen cried. "You can go and get that tray
        of yours!"}
}
{
}

\Verse{244}
{
    \zA{是巧宗儿，你们都得了，难道不许 我得一遭儿吗？”}
}
{
    \eA{"You should let me once go for something!" Ch'ing Wen pleaded. "Whenever any
        lucky chance has turned up, you've invariably grabbed it;}
}
{
}

\Verse{245}
{
    \zA{麝月笑道：“统共秋丫头得了一遭儿衣裳，那里今儿又巧，你也 遇见找衣裳不成？”}
}
{
    \eA{and can it be that you won't let me have a single turn?" "Altogether," She
        Yüeh said laughingly, "that girl Ch'iu Wen got a few clothes just once; can
        such a lucky coincidence present itself again today that you too should find
        them engaged in searching for clothes?"}
}
{
}

\Verse{246}
{
    \zA{晴雯冷笑道：“虽然碰不见衣裳，或者太太看见我勤谨，也把 太太的公费里一个月分出二两银子来给我，也定不得。”}
}
{
    \eA{"Albeit I mayn't come across any clothes," Ch'ing Wen rejoined with a
        sardonic smile, "our Madame Wang may notice how diligent I am, and apportion
        me a couple of taels out of her public expenses; there's no saying."
        Continuing,}
}
{
}

\Verse{247}
{
    \zA{说着，又笑道：“你们别 和我装神弄鬼的，什么事我不知道！”}
}
{
    \eA{"Don't you people," she laughed, "try and play your pranks with me; for is
        there anything that I don't twig?" As she spoke, she ran outside. Ch'iu Wen
        too left the room in her company;}
}
{
}

\Verse{248}
{
    \zA{一面说，一面往外跑了。}
}
{
    \eA{but she repaired to T'an Ch'un's quarters and fetched the tray.}
}
{
}

\Verse{249}
{
    \zA{秋纹也同他出来， 自去探春那里取了碟子来。}
}
{
    \eA{Hsi Jen then got everything ready. Calling an old nurse attached to the same
        place as herself, Sung by name, "Just go first and wash,}
}
{
}

\Verse{250}
{
    \zA{袭人打点齐备东西，叫过本处的一个老宋妈妈来，向他说道：“你去好生梳洗 了，换了出门的衣裳来，回来打发你给史大姑娘送东西去。”}
}
{
    \eA{comb your hair and put on your out-of-door clothes," she said to her, "and
        then come back as I want to send you at once with a present to Miss Shih."
        "Miss," urged the nurse Sung, "just give me what you have; and, if you have
        any message, tell it me; so that when I've tidied myself I may go
        straightway." Hsi Jen, at this proposal,}
}
{
}

\Verse{251}
{
    \zA{宋妈妈道：“姑娘只 管交给我，有话说与我，我收拾了就好一顺去。”}
}
{
    \eA{brought two small twisted wire boxes; and, opening first the one in which
        were two kinds of fresh fruits, consisting of caltrops and "chicken head"
        fruit, and afterwards uncovering the other,}
}
{
}

\Verse{252}
{
    \zA{袭人听说，便端过两个小摄丝盒 子来。}
}
{
    \eA{containing a tray with new cakes, made of chestnut powder, and steamed in
        sugar, scented with the olea,}
}
{
}

\Verse{253}
{
    \zA{先揭开一个，里面装的是红菱、鸡头两样鲜果；又揭开那个，是一碟子桂花 糖蒸的新栗粉糕。}
}
{
    \eA{"All these fresh fruits are newly plucked this year from our own garden,"
        she observed; "our Mr. Secundus sends them to Miss Shih to taste. The other
        day, too, she was quite taken with this cornelian tray so let her keep it
        for her use. In this silk bag she'll find the work,}
}
{
}

\Verse{254}
{
    \zA{又说道：“这都是今年咱们这里园里新结的果子，宝二爷送来给 姑娘尝尝。}
}
{
    \eA{which she asked me some time ago to do for her. (Tell her) that she mustn't
        despise it for its coarseness, but make the best of it and turn it to some
        account. Present respects to her from our part and inquire after her health
        on behalf of Mr. Pao-yü;}
}
{
}

\Verse{255}
{
    \zA{再前日姑娘说这玛瑙碟子好，姑娘就留下玩罢。}
}
{
    \eA{that will be all there's to say." "Has Mr. Pao, I wonder, anything more for
        me to tell her?" the nurse Sung added, "Miss, do go and inquire, so that on
        my return,}
}
{
}

\Verse{256}
{
    \zA{这绢包儿里头是姑娘前 日叫我做的活计，姑娘别嫌粗糙，将就着用罢。}
}
{
    \eA{he mayn't again say that I forgot." "He was just now," Hsi Jen consequently
        asked Ch'iu Wen, "over there in Miss Tertia's rooms, wasn't he?" "They were
        all assembled there,}
}
{
}

\Verse{257}
{
    \zA{替二爷问好，替我们请安，就是了。”}
}
{
    \eA{deliberating about starting some poetical society or other," Ch'iu Wen
        explained, "and they all wrote verses too.}
}
{
}

\Verse{258}
{
    \zA{宋妈妈道：“宝二爷不知还有什么说的？}
}
{
    \eA{But I fancy he's got no message to give you; so you might as well start."
        After this assurance, nurse Sung forthwith took the things,}
}
{
}

\Verse{259}
{
    \zA{姑娘再问问去，回来别又说忘了。”}
}
{
    \eA{and quitted the apartment. When she had changed her clothes and arranged her
        hair,}
}
{
}

\Verse{260}
{
    \zA{袭人 因问秋纹：“方才可是在三姑娘那里么？”}
}
{
    \eA{Hsi Jen further enjoined them to go by the back door, where there was a
        servant-boy, waiting with a curricle. Nurse Sung thereupon set out on her
        errand.}
}
{
}

\Verse{261}
{
    \zA{秋纹道：“他们都在那里商议起什么诗 社呢，又是做诗。}
}
{
    \eA{But we will leave her for the present. In a little time Pao-yü came back.
        After first cursorily glancing at the begonias for a time, he walked into
        his rooms,}
}
{
}

\Verse{262}
{
    \zA{想来没话，你只管去罢。”}
}
{
    \eA{and explained to Hsi Jen all about the poetical society they had managed to
        establish,}
}
{
}

\Verse{263}
{
    \zA{宋妈妈听了，便拿了东西出去，穿戴 了，袭人又嘱咐他：“你打后门去，有小子和车等着呢。”}
}
{
    \eA{Hsi Jen then told him that she had sent the nurse Sung along with some
        things, to Shih Hsiang-yün. As soon as Pao-yü heard this, he clapped his
        hands. "I forgot all about her!" he cried. "I knew very well that I had
        something to attend to; but I couldn't remember what it was!}
}
{
}

\Verse{264}
{
    \zA{宋妈妈去了，不在话下。}
}
{
    \eA{Luckily, you've alluded to her! I was just meaning to ask her to come,}
}
{
}

\Verse{265}
{
    \zA{一时宝玉回来，先忙着看了一回海棠，至屋里告诉袭人起诗社的事，袭人也把 打发宋妈妈给史湘云送东西去的话告诉了宝玉。}
}
{
    \eA{for what fun will there be in this poetical society without her?" "Is this
        of any serious import?" Hsi Jen reasoned with him. "It's all, for the mere
        sake of recreation! She's not however able to go about at her own free will
        as you people do. Nor can she at home have her own way.}
}
{
}

\Verse{266}
{
    \zA{宝玉听了，拍手道：“偏忘了他! 我只觉心里有件事，只是想不起来，亏你提起来，正要请他去。}
}
{
    \eA{When you therefore let her know, it won't again rest with her, however
        willing she may be to avail herself of your invitation. And if she can't
        come, she will long and crave to be with you all, so isn't it better that
        you shouldn't be the means of making her unhappy?"}
}
{
}

\Verse{267}
{
    \zA{这诗社里要少了他， 还有个什么意思！”}
}
{
    \eA{"Never mind!" responded Pao-yü. "I'll tell our venerable senior to despatch
        some one to bring her over." But in the middle of their conversation,}
}
{
}

\Verse{268}
{
    \zA{袭人劝道：“什么要紧，不过玩意儿。}
}
{
    \eA{nurse Sung returned already from her mission, and expressed to him,
        (Hsiang-yün's) acknowledgment; and to Hsi Jen her thanks for the trouble.}
}
{
}

\Verse{269}
{
    \zA{他比不得你们自在，家 里又作不得主儿。}
}
{
    \eA{"She also inquired," the nurse proceeded, "what you, master Secundus,}
}
{
}

\Verse{270}
{
    \zA{告诉他，他要来又由不得他，要不来他又牵肠挂肚的，没的叫他 不受用。”}
}
{
    \eA{were up to, and I told her that you had started some poetical club or other
        with the young ladies and that you were engaged in writing verses. Miss Shih
        wondered why it was, if you were writing verses, that you didn't even
        mention anything to her;}
}
{
}

\Verse{271}
{
    \zA{宝玉道：“不妨事，我回老太太，打发人接他去。”}
}
{
    \eA{and she was extremely distressed about it." Pao-yü, at these words, turned
        himself round and betook himself immediately into his grandmother's
        apartments,}
}
{
}

\Verse{272}
{
    \zA{正说着，宋妈妈已 经回来道生受，给袭人道乏，又说：“问二爷做什么呢，我说：‘和姑娘们起什么 诗社做诗呢。’}
}
{
    \eA{where he did all that lay in his power to urge her to depute servants to go
        and fetch her. "It's too late to-day," dowager lady Chia answered; "they'll
        go tomorrow, as soon as it's daylight." Pao-yü had no other course but to
        accede to her wishes. He, however, retraced his steps back to his room with
        a heavy heart.}
}
{
}

\Verse{273}
{
    \zA{史姑娘道，他们做诗，也不告诉他去。}
}
{
    \eA{On the morrow, at early dawn, he paid another visit to old lady Chia and
        brought pressure to bear on her until she sent some one for her.}
}
{
}

\Verse{274}
{
    \zA{急的了不得！”}
}
{
    \eA{Soon after midday, Shih Hsiang-yün arrived.}
}
{
}

\Verse{275}
{
    \zA{宝玉听了， 转身便往贾母处来，立逼着叫人接去。}
}
{
    \eA{Pao-yü felt at length much relieved in his mind. Upon meeting her, he
        recounted to her all that had taken place from beginning to end.}
}
{
}

\Verse{276}
{
    \zA{贾母因说：“今儿天晚了，明日一早去。”}
}
{
    \eA{His purpose was likewise to let her see the poetical composition, but Li Wan
        and the others remonstrated.}
}
{
}

\Verse{277}
{
    \zA{宝玉只得罢了。}
}
{
    \eA{"Don't," they said, "allow her to see them!}
}
{
}

\Verse{278}
{
    \zA{回来闷闷的，次日一早，便又往贾母处来催逼人接去。}
}
{
    \eA{First tell her the rhymes and number of feet; and, as she comes late, she
        should, as a first step, pay a penalty by conforming to the task we had to
        do. Should what she writes be good,}
}
{
}

\Verse{279}
{
    \zA{直到午后，湘云才来了，宝玉方放了心。}
}
{
    \eA{then she can readily be admitted as a member of the society; but if not
        good, she should be further punished by being made to stand a treat;}
}
{
}

\Verse{280}
{
    \zA{见面时，就把始末原由告诉他，又要 与他诗看。}
}
{
    \eA{after which, we can decide what's to be done." "You've forgotten to ask me
        round," Hsiang-yün laughed, "and I should,}
}
{
}

\Verse{281}
{
    \zA{李纨等因说道：“且别给他看，先说给他韵脚；他后来的，先罚他和了 诗。}
}
{
    \eA{after all, fine you people! But produce the metre; for though I don't excel
        in versifying, I shall exert myself to do the best I can, so as to get rid
        of every slur. If you will admit me into the club,}
}
{
}

\Verse{282}
{
    \zA{要好，就请入社；要不好，还要罚他一个东道儿再说。”}
}
{
    \eA{I shall be even willing to sweep the floors and burn the incense." When they
        all saw how full of fun she was, they felt more than ever delighted with her
        and they reproached themselves,}
}
{
}

\Verse{283}
{
    \zA{湘云笑道：“你们忘 了请我，我还要罚你们呢。}
}
{
    \eA{for having somehow or other managed to forget her on the previous day. But
        they lost no time in telling her the metre of the verses.}
}
{
}

\Verse{284}
{
    \zA{就拿韵来，我虽不能，只得勉强出丑。}
}
{
    \eA{Shih Hsiang-yün was inwardly in ecstasies. So much so, that she could not
        wait to beat the tattoo and effect any alterations.}
}
{
}

\Verse{285}
{
    \zA{容我入社，扫地 焚香，我也情愿。”}
}
{
    \eA{But having succeeded, while conversing with her cousins, in devising a
        stanza in her mind,}
}
{
}

\Verse{286}
{
    \zA{众人见他这般有趣，越发喜欢，都埋怨：“昨日怎么忘了他呢！”}
}
{
    \eA{she promptly inscribed it on the first piece of paper that came to hand. "I
        have," she remarked, with a precursory smile, "stuck to the metre and
        written two stanzas. Whether they be good or bad,}
}
{
}

\Verse{287}
{
    \zA{遂忙告诉他诗韵。}
}
{
    \eA{I cannot say; all I've kept in view was to simply comply with your wishes."}
}
{
}

\Verse{288}
{
    \zA{湘云一心兴头，等不得推敲删改，一面只管和人说着话，心内早已和成，即用 随便的纸笔录出，先笑说道：“我却依韵和了两首，好歹我都不知，不过应命而已。”}
}
{
    \eA{So speaking, she handed her paper to the company. "We thought our four
        stanzas," they observed, "had so thoroughly exhausted everything that could
        be imagined on the subject that another stanza was out of the question, and
        there you've devised a couple more! How could there be so much to say? These
        must be mere repetitions of our own sentiments."}
}
{
}

\Verse{289}
{
    \zA{说着，递与众人。}
}
{
    \eA{While bandying words, they perused her two stanzas.}
}
{
}

\Verse{290}
{
    \zA{众人道：“我们四首也算想绝了，再一首也不能了，你倒弄了两 首!那里有许多话说？}
}
{
    \eA{They found this to be their burden: No. 1. The fairies yesterday came down
        within the city gates, And like those gems, sown in the grassy field,
        planted one pot. How clear it is that the goddess of frost is fond of cold!
        It is no question of a pretty girl bent upon death!}
}
{
}

\Verse{291}
{
    \zA{必要重了我们的。”}
}
{
    \eA{Where does the snow, which comes in gloomy weather,}
}
{
}

\Verse{292}
{
    \zA{一面说，一面看时，只见那两首诗写道： 白海棠和韵 神仙昨日降都门，种得蓝田玉一盆。}
}
{
    \eA{issue from? The drops of rain increase the prints, left from the previous
        night. How the flowers rejoice that bards are not weary of song! But are
        they ever left to spend in peace a day or night? No. 2. The "heng chih"
        covered steps lead to the creeper-laden door.}
}
{
}

\Verse{293}
{
    \zA{自是霜娥偏爱冷，非关倩女欲离魂。}
}
{
    \eA{How fit to plant by the corner of walls; how fit for pots? The flowers so
        relish purity that they can't find a mate.}
}
{
}

\Verse{294}
{
    \zA{秋阴捧出何方雪？}
}
{
    \eA{Easy in autumn snaps the soul of sorrow-wasted man.}
}
{
}

\Verse{295}
{
    \zA{雨渍添来隔宿痕。}
}
{
    \eA{The tears, which from the jade-like candle drip,}
}
{
}

\Verse{296}
{
    \zA{却喜诗人吟不倦，肯令寂寞度朝昏？}
}
{
    \eA{dry in the wind. The crystal-like portiere asunder rends Selene's rays.
        Their private feelings to the moon goddess they longed to tell,}
}
{
}

\Verse{297}
{
    \zA{其　二 蘅芷阶通萝薜门，也宜墙角也宜盆。}
}
{
    \eA{But gone, alas! is the lustre she shed on the empty court! Every line filled
        them with wonder and admiration. What they read, they praised.}
}
{
}

\Verse{298}
{
    \zA{花因喜洁难寻偶，人为悲秋易断魂。}
}
{
    \eA{"This," they exclaimed, with one consent, "is not writing verses on the
        begonia for no purpose! We must really start a Begonia Society!"}
}
{
}

\Verse{299}
{
    \zA{玉烛滴干风里泪，晶帘隔破月中痕。}
}
{
    \eA{"To-morrow," Shih Hsiang-yün proposed, "first fine me by making me stand a
        treat, and letting me be the first to convene a meeting;}
}
{
}

\Verse{300}
{
    \zA{幽情欲向嫦娥诉，无那虚廊月色昏。}
}
{
    \eA{may I?" "This would be far better!" they all assented. So producing also the
        verses, composed the previous day,}
}
{
}

\Verse{301}
{
    \zA{众人看一句惊讶一句，看到了赞到了都说：“这个不枉做了海棠诗!真该要起‘海 棠社’了。”}
}
{
    \eA{they submitted them to her for criticism. In the evening, Hsiang-yün came at
        the invitation of Pao-ch'ai, to the Heng Wu Yüan to put up with her for the
        night. By lamplight, Hsiang-yün consulted with her how she was to play the
        hostess and fix upon the themes;}
}
{
}

\Verse{302}
{
    \zA{湘云道：“明日先罚我个东道儿，就让我先邀一社，可使得？”}
}
{
    \eA{but, after lending a patient ear to all her proposals for a long time,
        Pao-ch'ai thought them so unsuitable for the occasion, that turning towards
        her,}
}
{
}

\Verse{303}
{
    \zA{众人 道：“这更妙了。”}
}
{
    \eA{she raised objections. "If you want," she said, "to hold a meeting,}
}
{
}

\Verse{304}
{
    \zA{因又将昨日的诗与他评论了一回。}
}
{
    \eA{you have to pay the piper. And albeit it's for mere fun, you have to make
        every possible provision;}
}
{
}

\Verse{305}
{
    \zA{至晚，宝钗将湘云邀往蘅芜院去安歇。}
}
{
    \eA{for while consulting your own interests, you must guard against giving
        umbrage to people.}
}
{
}

\Verse{306}
{
    \zA{湘云灯下计议如何设东拟题。}
}
{
    \eA{In that case every one will afterwards be happy and contented. You count for
        nothing too in your own home;}
}
{
}

\Verse{307}
{
    \zA{宝钗听他 说了半日，皆不妥当，因向他说道：“既开社，就要作东。}
}
{
    \eA{and the whole lump sum of those few tiaos, you draw each month, are not
        sufficient for your own wants, and do you now also wish to burden yourself
        with this useless sort of thing?}
}
{
}

\Verse{308}
{
    \zA{虽然是个玩意儿，也要 瞻前顾后；又要自己便宜，又要不得罪了人，然后方大家有趣。}
}
{
    \eA{Why, if your aunt gets wind of it, won't she be more incensed with you than
        ever! What's more, even though you might fork out all the money you can call
        your own to bear the outlay of this entertainment with, it won't be anything
        like enough,}
}
{
}

\Verse{309}
{
    \zA{你家里你又做不得 主，一个月统共那几吊钱，你还不够使。}
}
{
    \eA{and can it possibly be, pray, that you would go home for the express purpose
        of requisitioning the necessary funds? Or will you perchance ask for some
        from in here?"}
}
{
}

\Verse{310}
{
    \zA{这会子又干这没要紧的事，你婶娘听见了 越发抱怨你了。}
}
{
    \eA{This long tirade had the effect of bringing the true facts of the case to
        Hsiang-yün's notice, and she began to waver in a state of uncertainty.}
}
{
}

\Verse{311}
{
    \zA{况且你就都拿出来，做这个东也不够，难道为这个家去要不成？}
}
{
    \eA{"I have already fixed upon a plan in my mind," Pao-ch'ai resumed. "There's
        an assistant in our pawnshop from whose family farm come some splendid
        crabs.}
}
{
}

\Verse{312}
{
    \zA{还 是和这里要呢？”}
}
{
    \eA{Some time back, he sent us a few as a present,}
}
{
}

\Verse{313}
{
    \zA{一席话提醒了湘云，倒踌蹰起来。}
}
{
    \eA{and now, starting from our venerable senior and including the inmates of the
        upper quarters,}
}
{
}

\Verse{314}
{
    \zA{宝钗道：“这个我已经有个主 意了。}
}
{
    \eA{most of them are quite in love with crabs. It was only the other day that my
        mother mentioned that she intended inviting our worthy ancestor into the
        garden to look at the olea flowers and partake of crabs,}
}
{
}

\Verse{315}
{
    \zA{我们当铺里有个伙计，他们地里出的好螃蟹，前儿送了几个来。}
}
{
    \eA{but she has had her hands so full that she hasn't as yet asked her round. So
        just you now drop the poetical meeting, and invite the whole crowd to a
        show; and if we wait until they go,}
}
{
}

\Verse{316}
{
    \zA{现在这里的 人，从老太太起，连上屋里的人，有多一半都是爱吃螃蟹的，前日姨娘还说要请老 太太在园里赏桂花、吃螃蟹，因为有事，还没有请。}
}
{
    \eA{won't we be able to indite as many poems as we like? But let me speak to my
        brother and ask him to let us have several baskets of the fattest and
        largest crabs he can get, and to also go to some shop and fetch several jars
        of luscious wine. And if we then lay out four or five tables with plates
        full of refreshments, won't we save trouble and all have a jolly time as
        well?" As soon as Hsiang-yün heard (the alternative proposed by Pao-ch'ai,)}
}
{
}

\Verse{317}
{
    \zA{你如今且把诗社别提起，只普 同一请，等他们散了，咱们有多少诗做不得的？}
}
{
    \eA{she felt her heart throb with gratitude and in most profuse terms she
        praised her for her forethought. "The proposal I've made." Pao-ch'ai pursued
        smilingly; "is prompted entirely by my sincere feelings for you;}
}
{
}

\Verse{318}
{
    \zA{我和我哥哥说，要他几篓极肥极大 的螃蟹来，再往铺子里取上几坛好酒来，再备四五桌果碟子，岂不又省事，又大家 热闹呢？”}
}
{
    \eA{so whatever you do don't be touchy and imagine that I look down upon you;
        for in that case we two will have been good friends all in vain. But if you
        won't give way to suspicion, I'll be able to tell them at once to go and get
        things ready." "My dear cousin," eagerly rejoined Hsiang-yün, a smile on her
        lips, "if you say these things it's you who treat me with suspicion;}
}
{
}

\Verse{319}
{
    \zA{湘云听了，心中自是感服，极赞想的周到。}
}
{
    \eA{for no matter how foolish a person I may be, as not to even know what's good
        and bad, I'm still a human being!}
}
{
}

\Verse{320}
{
    \zA{宝钗又笑道：“我是一片真 心为你的话，你可别多心，想着我小看了你，咱们两个就白好了。}
}
{
    \eA{Did I not regard you, cousin, in the same light as my own very sister, I
        wouldn't last time have had any wish or inclination to disclose to you every
        bit of those troubles, which ordinarily fall to my share at home."}
}
{
}

\Verse{321}
{
    \zA{你要不多心，我 就好叫他们办去。”}
}
{
    \eA{After listening to these assurances, Pao-ch'ai summoned a matron and bade
        her go out and tell her master,}
}
{
}

\Verse{322}
{
    \zA{湘云忙笑道：“好姐姐!你这么说，倒不是真心待我了。}
}
{
    \eA{Hsüeh P'an, to procure a few hampers of crabs of the same kind as those
        which were sent on the previous occasion. "Our venerable senior," (she
        said,)}
}
{
}

\Verse{323}
{
    \zA{我凭 怎么胡涂，连个好歹也不知，还是个人吗!我要不把姐姐当亲姐姐待，上回那些家 常烦难事，我也不肯尽情告诉你了。”}
}
{
    \eA{"and aunt Wang are asked to come to-morrow after their meal and admire the
        olea flowers, so mind, impress upon your master to please not forget, as
        I've already to-day issued the invitations." The matron walked out of the
        garden and distinctly delivered the message.}
}
{
}

\Verse{324}
{
    \zA{宝钗听说，便唤一个婆子来：“出去和大爷 说，照前日的大螃蟹要几篓来，明日饭后请老太太、姨娘赏桂花。}
}
{
    \eA{But, on her return, she brought no reply. During this while, Pao-ch'ai
        continued her conversation with Hsiang-yün. "The themes for the verses," she
        advised her, "mustn't also be too out-of-the-way. Just search the works of
        old writers,}
}
{
}

\Verse{325}
{
    \zA{你说与大爷：好 歹别忘了，我今儿已经请下人了。”}
}
{
    \eA{and where will you find any eccentric and peculiar subjects, or any extra
        difficult metre! If the subject be too much out-of-the-way and the metre too
        difficult,}
}
{
}

\Verse{326}
{
    \zA{那婆子出去说明，回来无话。}
}
{
    \eA{one cannot get good verses. In a word, we are a mean lot and our verses are
        certain,}
}
{
}

\Verse{327}
{
    \zA{这里宝钗又向湘云道：“诗题也别过于新巧了，你看古人中那里有那些刁钻古 怪的题目和那极险的韵呢？}
}
{
    \eA{I fear, to consist of mere repetitions. Nor is it advisable for us to aim at
        excessive originality. The first thing for us to do is to have our ideas
        clear, as our language will then not be commonplace.}
}
{
}

\Verse{328}
{
    \zA{若题目过于新巧，韵过于险，再不得好诗，倒小家子气。}
}
{
    \eA{In fact, this sort of thing is no vital matter; spinning and needlework are,
        in a word, the legitimate duties of you and me. Yet, if we can at any time
        afford the leisure,}
}
{
}

\Verse{329}
{
    \zA{诗固然怕说熟话，然也不可过于求生；头一件，只要主意清新，措词就不俗了。}
}
{
    \eA{it's only right and proper that we should take some book, that will benefit
        both body and mind, and read a few chapters out of it." Hsiang-yün simply
        signified her assent. "I'm now cogitating in my mind,"}
}
{
}

\Verse{330}
{
    \zA{究 竟这也算不得什么，还是纺绩针黹是你我的本等。}
}
{
    \eA{she then laughingly remarked, "that as the verses we wrote yesterday treated
        of begonias, we should, I think, compose on this occasion some on
        chrysanthemums,}
}
{
}

\Verse{331}
{
    \zA{一时闲了，倒是把那于身心有益 的书看几章，却还是正经。”}
}
{
    \eA{eh? What do you say?" "Chrysanthemums are in season," Pao-ch'ai replied.
        "The only objection to them is that too many writers of old have made them
        the subject of their poems."}
}
{
}

\Verse{332}
{
    \zA{湘云只答应着，因笑道：“我心里想着，昨日做了海 棠诗，我如今要做个菊花诗如何？”}
}
{
    \eA{"I also think so," Hsiang-yün added, "so that, I fear, we shall only be
        following in their footsteps." After some reflection, Pao-ch'ai exclaimed,
        "I've hit upon something! If we take, for the present instance,}
}
{
}

\Verse{333}
{
    \zA{宝钗道：“菊花倒也合景，只是前人太多了。”}
}
{
    \eA{the chrysanthemums as a secondary term, and man as the primary, we can,
        after all, select several themes. But they must all consist of two
        characters:}
}
{
}

\Verse{334}
{
    \zA{湘云道：“我也是这么想着，恐怕落套。”}
}
{
    \eA{the one, an empty word; the other, a full one. The full word might be
        chrysanthemums; while for the empty one,}
}
{
}

\Verse{335}
{
    \zA{宝钗想了一想，说道：“有了。}
}
{
    \eA{we might employ some word in general use. In this manner, we shall, on one
        hand, sing the chrysanthemum;}
}
{
}

\Verse{336}
{
    \zA{如今以 菊花为宾，以人为主，竟拟出几个题目来，都要两个字，一个虚字一个实字。}
}
{
    \eA{and, on the other, compose verses on the theme. And as old writers have not
        written much in this style, it will be impossible for us to drift into the
        groove of their ideas. Thus in versifying on the scenery and in singing the
        objects,}
}
{
}

\Verse{337}
{
    \zA{实字 就用‘菊’字，虚字便用通用门的。}
}
{
    \eA{we will, in both respects, combine originality with liberality of thought."
        "This is all very well,"}
}
{
}

\Verse{338}
{
    \zA{如此，又是咏菊，又是赋事，前人虽有这么做 的，还不很落套。}
}
{
    \eA{smiled Hsiang-yün. "The only thing is what kind of empty words will, I
        wonder, be best to use? Just you first think of one and let me see."
        Pao-ch'ai plunged in thought for a time,}
}
{
}

\Verse{339}
{
    \zA{赋景咏物两关着，也倒新鲜大方。”}
}
{
    \eA{after which she laughingly remarked: "Dream of chrysanthemums is good."
        "It's positively good!" Hsiang-yün smiled.}
}
{
}

\Verse{340}
{
    \zA{湘云笑道：“很好，只是不 知用什么虚字才好？}
}
{
    \eA{"I've also got one: 'the Chrysanthemum shadow,' will that do?" "Well
        enough," Pao-ch'ai answered, "the only objection is that people have written
        on it;}
}
{
}

\Verse{341}
{
    \zA{你先想一个我听听。”}
}
{
    \eA{yet if the themes are to be many, we might throw this in. I've got another
        one too!"}
}
{
}

\Verse{342}
{
    \zA{宝钗想了一想，笑道：“‘菊梦’就好。”}
}
{
    \eA{"Be quick, and tell it!" Hsiang-yün urged. "What do you say to 'ask the
        Chrysanthemums?'"}
}
{
}

\Verse{343}
{
    \zA{湘云笑道：“果然好。}
}
{
    \eA{Pao-ch'ai observed. Hsiang-yün clapped her hand on the table.}
}
{
}

\Verse{344}
{
    \zA{我也有一个： ‘菊影’可使得？”}
}
{
    \eA{"Capital," she cried. "I've thought of one also." She then quickly
        continued,}
}
{
}

\Verse{345}
{
    \zA{宝钗道：“也罢了，只是也有人做过。}
}
{
    \eA{"It is, search for chrysanthemums; what's your idea about it?" Pao-ch'ai
        thought that too would do very well.}
}
{
}

\Verse{346}
{
    \zA{若题目多，这个也搭的 上。}
}
{
    \eA{"Let's choose ten of them first," she next proposed; "and afterwards note
        them down!"}
}
{
}

\Verse{347}
{
    \zA{我又有了一个。”}
}
{
    \eA{While talking, they rubbed the ink and moistened the pens.}
}
{
}

\Verse{348}
{
    \zA{湘云道：“快说出来。”}
}
{
    \eA{These preparations over, Hsiang-yün began to write,}
}
{
}

\Verse{349}
{
    \zA{宝钗道：“‘问菊’如何？”}
}
{
    \eA{while Pao-ch'ai enumerated the themes. In a short time, they got ten of
        them.}
}
{
}

\Verse{350}
{
    \zA{湘云 拍案叫妙，因接说道：“我也有了：‘访菊’好不好？”}
}
{
    \eA{"Ten don't form a set," Hsiang-yün went on to smilingly suggest, after
        reading them over. "We'd better complete them by raising their number to
        twelve;}
}
{
}

\Verse{351}
{
    \zA{宝钗也赞有趣。}
}
{
    \eA{they'll then also be on the same footing as people's pictures and books."}
}
{
}

\Verse{352}
{
    \zA{因说道： “索性拟出十个来，写上再来。”}
}
{
    \eA{Hearing this proposal, Pao-ch'ai devised another couple of themes, thus
        bringing them to a dozen.}
}
{
}

\Verse{353}
{
    \zA{说着，二人研墨蘸笔，湘云便写，宝钗便念，一 时凑了十个。}
}
{
    \eA{"Well, since we've got so far," she pursued, "let's go one step further and
        copy them out in their proper order, putting those that are first, first;
        and those that come last,}
}
{
}

\Verse{354}
{
    \zA{湘云看了一遍，又笑道：“十个还不成幅，索性凑成十二个，就全了， 也和人家的字画册页一样。”}
}
{
    \eA{last." "It would be still better like that," Hsiang-yün acquiesced, "as
        we'll be able to make up a 'chrysanthemum book.'" "The first stanza should
        be: 'Longing for chrysanthemums,'"}
}
{
}

\Verse{355}
{
    \zA{宝钗听说，又想了两个，一共凑成十二个，说道：“既 这么着，一发编出个次序来。”}
}
{
    \eA{Pao-chai said, "and as one cannot get them by wishing, and has, in
        consequence, to search for them, the second should be 'searching for
        chrysanthemums.' After due search, one finds them, and plants them, so the
        third must be:}
}
{
}

\Verse{356}
{
    \zA{湘云道：“更妙，竟弄成个菊谱了。”}
}
{
    \eA{'planting chrysanthemums.' After they've been planted, they, blossom, and
        one faces them and enjoys them,}
}
{
}

\Verse{357}
{
    \zA{宝钗道：“起首是《忆菊》；忆之不得，故访，第二是《访菊》。}
}
{
    \eA{so the fourth should be 'facing the chrysanthemums.' By facing them, one
        derives such excessive delight that one plucks them and brings them in and
        puts them in vases for one's own delectation,}
}
{
}

\Verse{358}
{
    \zA{访之既得， 便种，第三是《种菊》。}
}
{
    \eA{so the fifth must be 'placing chrysanthemums in vases.' If no verses are
        sung in their praise,}
}
{
}

\Verse{359}
{
    \zA{种既盛开，故相对而赏，第四是《对菊》。}
}
{
    \eA{after they've been placed in vases, it's tantamount to seeing no point of
        beauty in chrysanthemums, so the sixth must be 'sing about chrysanthemums.'}
}
{
}

\Verse{360}
{
    \zA{相对而兴有馀， 故折来供瓶为玩，第五是《供菊》。}
}
{
    \eA{After making them the burden of one's song, one can't help representing them
        in pictures. The seventh place should therefore be conceded to 'drawing
        chrysanthemums.'}
}
{
}

\Verse{361}
{
    \zA{既供而不吟，亦觉菊无彩色，第六便是《咏菊》。}
}
{
    \eA{Seeing that in spite of all the labour bestowed on the drawing of
        chrysanthemums, the fine traits there may be about them are not yet,}
}
{
}

\Verse{362}
{
    \zA{既入词章，不可以不供笔墨，第七便是《画菊》。}
}
{
    \eA{in fact, apparent, one impulsively tries to find them out by inquiries, so
        the eighth should be 'asking the chrysanthemums.'}
}
{
}

\Verse{363}
{
    \zA{既然画菊，若是默默无言，究竟 不知菊有何妙处，不禁有所问，第八便是《问菊》。}
}
{
    \eA{As any perception, which the chrysanthemums might display in fathoming the
        questions set would help to make the inquirer immoderately happy, the ninth
        must be 'pinning the chrysanthemums in the hair.' And as after everything
        has been accomplished,}
}
{
}

\Verse{364}
{
    \zA{菊若能解语，使人狂喜不禁， 便越要亲近他，第九竟是《簪菊》。}
}
{
    \eA{that comes within the sphere of man, there will remain still some
        chrysanthemums about which something could be written, two stanzas on the
        'shadow of the chrysanthemums,'}
}
{
}

\Verse{365}
{
    \zA{如此人事虽尽，犹有菊之可咏者，《菊影》《菊 梦》二首，续在第十、第十一。}
}
{
    \eA{and the 'dream about chrysanthemums' must be tagged on as numbers ten and
        eleven. While the last section should be 'the withering of the
        chrysanthemums' so as to bring to a close the sentiments expressed in the
        foregoing subjects.}
}
{
}

\Verse{366}
{
    \zA{末卷便以《残菊》总收前题之感。}
}
{
    \eA{In this wise the fine scenery and fine doings of the third part of autumn,
        will both alike be included in our themes."}
}
{
}

\Verse{367}
{
    \zA{这便是三秋的妙 景妙事都有了。”}
}
{
    \eA{Hsiang-yün signified her approval, and taking the list she copied it out
        clean.}
}
{
}

\Verse{368}
{
    \zA{湘云依言将题录出，又看了一回，又问：“该限何韵？”}
}
{
    \eA{But after once more passing her eye over it, she went on to inquire what
        rhymes should be determined upon. "I do not, as a rule, like hard-and-fast
        rhymes,"}
}
{
}

\Verse{369}
{
    \zA{宝钗道： “我平生最不喜限韵，分明有好诗，何苦为韵所缚？}
}
{
    \eA{Pao-ch'ai retorted. "It's evident enough that we can have good verses
        without them, so what's the use of any rhymes to shackle us? Don't let us
        imitate that mean lot of people.}
}
{
}

\Verse{370}
{
    \zA{咱们别学那小家派。}
}
{
    \eA{Let's simply choose our subject and pay no notice to rhymes.}
}
{
}

\Verse{371}
{
    \zA{只出题， 不拘韵：原为大家偶得了好句取乐，并不为以此难人。”}
}
{
    \eA{Our main object is to see whether we cannot by chance hit upon some
        well-written lines for the sake of fun. It isn't to make this the means of
        subjecting people to perplexities."}
}
{
}

\Verse{372}
{
    \zA{湘云道：“这话很是。}
}
{
    \eA{"What you say is perfectly right,"}
}
{
}

\Verse{373}
{
    \zA{既 这样，自然大家的诗还进一层。}
}
{
    \eA{Hsiang-yün observed. "In this manner our poetical composition will improve
        one step higher.}
}
{
}

\Verse{374}
{
    \zA{但只咱们五个人，这十二个题目，难道每人作十二 首不成？”}
}
{
    \eA{But we only muster five members, and there are here twelve themes. Is it
        likely that each one of us will have to indite verses on all twelve?" "That
        would be far too hard on the members!"}
}
{
}

\Verse{375}
{
    \zA{宝钗道：“那也太难人了。}
}
{
    \eA{Pao-ch'ai rejoined. "But let's copy out the themes clean,}
}
{
}

\Verse{376}
{
    \zA{将这题目誊好，都要七言律诗，明日贴在墙 上，他们看了，谁能那一个就做那一个。}
}
{
    \eA{for lines with seven words will have to be written on every one, and stick
        them to-morrow on the wall for general perusal. Each member can write on the
        subject which may be most in his or her line. Those, with any ability,}
}
{
}

\Verse{377}
{
    \zA{有力量者十二首都做也可，不能的作一首 也可，高才捷足者为尊。}
}
{
    \eA{may choose all twelve. While those, with none, may only limit themselves to
        one stanza. Both will do. Those, however, who will show high mental
        capacity, combined with quickness,}
}
{
}

\Verse{378}
{
    \zA{若十二首已全，便不许他赶着又做，罚他便完了。”}
}
{
    \eA{will be held the best. But any one, who shall have completed all twelve
        themes, won't be permitted to hasten and begin over again;}
}
{
}

\Verse{379}
{
    \zA{湘云 道：“这也罢了。”}
}
{
    \eA{we'll have to fine such a one, and finish." "Yes, that will do,"}
}
{
}

\Verse{380}
{
    \zA{二人商议妥贴，方才息灯安寝。}
}
{
    \eA{assented Hsiang-yün. But after settling everything satisfactorily, they
        extinguished the lamp and went to bed.}
}
{
}

\Verse{381}
{
    \zA{要知端底，下回分解。}
}
{
    \eA{Reader, do you want to know what subsequently took place?}
}
{
}

\Verse{382}
{
    \zA{(本章完)}
}
{
    \eA{If you do, then listen to what is contained in the way of explanation in the
        following chapter.}
}
{
}
